% Open Logic Project
%
% driver file open-logic-sample.tex to produce text on letter-size paper
% with standard layout and margins

% We use the memoir class for maximal flexibility of layout, but any
% class will do

\documentclass[b5paper]{memoir}

% \olpath has to point to the location of the OLP main
% directory/folder.  We're compiling from subdirectory courses/sample,
% so the main directory is two levels up.
\newcommand{\olpath}{../../}

% load all the Open Logic definitions. This will also load the
% local definitions in open-logic-sample-config.sty
% OpenLogic Project
% Master Style Package

% load various required standard packages
% let LaTeX find those

\RequirePackage[english]{babel}
\RequirePackage{amsmath}
\RequirePackage{amssymb}
\RequirePackage{extarrows}
\RequirePackage{mathdots}
\RequirePackage[OMLmathsfit]{isomath}
\RequirePackage{xparse}
\RequirePackage{mfirstuc}
\RequirePackage{subfiles}
\RequirePackage{framed}
\RequirePackage{graphicx}
\RequirePackage{multirow}
\RequirePackage{longtable}

\RequirePackage{natbib}
\setcitestyle{authoryear}

\RequirePackage{tikz}
\RequirePackage{tikzscale}
\usetikzlibrary{arrows,automata,intersections,positioning}

\RequirePackage{bussproofs}
\RequirePackage[tableaux]{prooftrees}

\newenvironment{oltableau}{\center\tableau{}} %wff format={anchor = base west}}}
               {\endtableau\endcenter}


\RequirePackage{wrapfig}
\RequirePackage{catchfile}
\RequirePackage{tocloft}
% Increase the space reserved for the section numbers in the table of contents
% (else some entries will look like "10.11Soundness" without any separation
% between the number and the title)
\setlength\cftsectionnumwidth{2.7em}
\setlength\cftpartnumwidth{2.7em}

\RequirePackage{xcolor}
\definecolor{dark-gray}{gray}{0.2}
\definecolor{reflex-blue}{RGB}{12, 28, 140}

\RequirePackage{hyperref}
\hypersetup{breaklinks,
  bookmarks,
  bookmarksopen,
  bookmarksopenlevel=1,
  colorlinks,
  urlcolor=dark-gray}
  %linkcolor=reflex-blue}

\RequirePackage{amsthm}
\RequirePackage{thmtools}
\RequirePackage[nameinlink]{cleveref} % always after hyperref!
\RequirePackage{nicefrac}

%\usepackage{pgf}

\makeatletter

% load all tool packages
% located in \olpath/sty/

% OpenLogic Project
% Package for crossreferencing

\NeedsTeXFormat{LaTeX2e}

\ProvidesFile{open-logic-referencing.sty}

\newcommand*{\theolpart}{udf}
\newcommand*{\theolchapter}{udf}
\newcommand*{\theolsection}{udf}
\newcommand*{\olfileid}[3]{
  \renewcommand*{\theolpart}{#1}
  \renewcommand*{\theolchapter}{#2}
  \renewcommand*{\theolsection}{#3}}
\newcommand*{\olchapterid}[2]{
  \renewcommand*{\theolpart}{#1}
  \renewcommand*{\theolchapter}{#2}}
\newcommand*{\olpartid}[1]{
  \renewcommand*{\theolpart}{#1}}

\newcommand*{\ollabel}[1]{%
  \label{\theolpart:\theolchapter:\theolsection:#1}}

\DeclareDocumentCommand \olref { o o o m } {%
    \cref{%
      \IfNoValueTF {#1} 
        {\theolpart:\theolchapter:\theolsection}
        {\IfNoValueTF {#2} 
          {\theolpart:\theolchapter:#1}
          {\IfNoValueTF {#3} 
            {\theolpart:#1:#2}
            {#1:#2:#3}}}:#4}}

\DeclareDocumentCommand \Olref { o o o m } {%
    \Cref{%
      \IfNoValueTF {#1} 
        {\theolpart:\theolchapter:\theolsection}
        {\IfNoValueTF {#2} 
          {\theolpart:\theolchapter:#1}
          {\IfNoValueTF {#3} 
            {\theolpart:#1:#2}
            {#1:#2:#3}}}:#4}}

% \olsection - every OLP text starts with this, starts a \section
% and sets a label. Refr to the section using \olref[part][chapter][section]{sec}

\DeclareDocumentCommand \olsection { o m } {%
  \IfNoValueTF {#1}{\csuse{\ol@sectioncs}{#2}}{\csuse{\ol@sectioncs}[#1]{#2}}%
  \ollabel{sec}}

\def\ol@sectioncs{section}

% \nosection - sectioning command that supresses section header.

\DeclareDocumentCommand \nosection { o m }{}

% \olchapter - make a chapter header and set a label for the chapter
% \olchapter[alt title]{part}{chapter}{title}
% Refer to the chapter using \olref[part][chapter][]{chap}

\DeclareDocumentCommand \olchapter { o m m m } {%
  \IfNoValueTF {#1}{\csuse{\ol@chaptercs}{#4}}{\csuse{\ol@chaptercs}[#1]{#4}}%
  \label{#2:#3::chap}}

\def\ol@defaultchaptercs{chapter}
\def\olsetchapter#1{\def\ol@chaptercs{#1}}
\def\olresetchapter{\let\ol@chaptercs\ol@defaultchaptercs}
\olresetchapter

% \olpart - make a part header and set a label for the part
% \olchapter[alt title]{part}{title}
% Refer to the part using \olref[part][][]{part}

\DeclareDocumentCommand \olpart { o m m } {%
  \IfNoValueTF {#1}{\csuse{\ol@partcs}{#3}}{\csuse{\ol@partcs}[#1]{#3}}%
  \label{#2:::part}}

\def\ol@partcs{part}
  
% \oliflabeldef - selectively include text iff a section label is
% defined

\DeclareDocumentCommand \oliflabeldef { m m m}{%
  \@ifundefined{r@#1}{#3}{#2}%
}

\crefformat{enumi}{#2(#1)#3}

% OpenLogic Project
% formulas

\NeedsTeXFormat{LaTeX2e}

\ProvidesFile{open-logic-formulas}


\RequirePackage{amsfonts}
\RequirePackage{amssymb}
\RequirePackage[OMLmathsfit]{isomath}

\def\fac{\text{!}}

% make ! active in math mode for formula meta variables

{\catcode`!=\active
\gdef!{\formula}
\let\formula\relax % default: do nothing
}
\mathcode`\!="8000 % set ! active in math mode

% latinformulas -- switch to latin formulas A, B, C, etc.

\def\ollatinformulas{\let\formula\relax}

% greekformulas -- formula metavariables are phi, psi, etc.
% Provides !A, !B etc in math mode to produce \phi, \psi etc.

\def\olgreekformulas
{\gdef\formula{\lookup{A/\varphi,B/\psi,C/\chi,D/\theta,E/\alpha,F/\beta,G/\gamma,H/\delta,K/\xi,L/\zeta,O/\omega,R/\rho,S/\sigma,T/\tau}}}

\def\olalphagreekformulas
{\gdef\formula{\lookup{A/\alpha,B/\beta,C/\gamma,D/\delta,E/\phi,F/\psi,G/\chi,H/\theta,K/\xi,L/\zeta,O/\omega,R/\rho,S/\sigma,T/\tau}}}

% args -- format a list of arguments separated by commas without commas

\newcommand*{\args}[1]{%
  \let\@argsep\@argsepinit
  \@for\@arg:=#1\do{%
    \@argsep\@arg}%
}

\def\@argsepinit{\let\@argsep\argsep}

\newcommand{\argsep}{}

% lookup#1#2 -- look up #2 in #1, a comma-separated list of pairs

\def\@match#1/#2{\def\match@input{#1}\def\match@result{#2}}
\newcommand*{\lookup}[2]{%
  \def\lookup@input{#2}
  \def\lookup@result{#2}%
  \def\do##1{\@match##1\if\lookup@input\match@input
    \let\lookup@result\match@result\fi}%
  \docsvlist{#1}%
  \lookup@result}

% OpenLogic Project
% Package for tokenizing text

\NeedsTeXFormat{LaTeX2e}
\ProvidesFile{open-logic-tokenize}

\RequirePackage{mfirstuc}
\RequirePackage{xparse}

% Tokens 
% ====== 

% Text that should be uniform and changeable Has four forms: in-text
% singular, in-text plural, initial singular, initial plural. Initial
% forms are usually capitalized. They are optional; if left out, we'll
% take the in-text version and make the first letter uppercase. 
% Optional star indicates nouns that take "an" rather than "a" as article.

% `\definetoken{switch}{token}{replacement text}` - defines a macro
% `\TOK@switch@token` to produce <replacement text>. `\definetokenuc`
% takes the first character and makes it uppercase.

\newcommand*{\definetoken}[3]{\expandafter\edef\csname TOK@#1@#2\endcsname{#3}}

\newcommand*{\definetokenuc}[3]{\definetokenucaux{#1}{#2}#3\relax}
\def\definetokenucaux #1#2#3#4\relax{
\uppercase{\gdef\firstcharuc{#3}}
\definetoken{#1}{#2}{\firstcharuc #4}}

% `\usetoken{switch}{token}` - produces the replacement text for
% switch/token, i.e., whatever the macro `\TOK@switch@token` was
% defined as

\newcommand*{\usetoken}[2]{\csname TOK@#1@#2\endcsname}

\DeclareDocumentCommand \settexttoken { m s m m o o } {%
    \definetoken{s}{#1}{#3}
    \definetoken{p}{#1}{#4}
    \IfBooleanTF {#2} 
        { \definetoken{a}{#1}{an} 
          \definetoken{A}{#1}{An}}
        { \definetoken{a}{#1}{a} 
          \definetoken{A}{#1}{A}}
    \IfNoValueTF {#5} 
        { \definetokenuc{S}{#1}{#3}}
        { \definetoken{S}{#1}{#5} }
    \IfNoValueTF {#6} 
        { \definetokenuc{P}{#1}{#4}}
        { \definetoken{P}{#1}{#6}}}

% `\@printtoken` - non-expandable wrapper macro for tokens, essentially
% what `!!` does: if first arg is `^`, we want uppercase, if next arg
% is `a` we want an article, if `{token}` followed by `s` we want
% plural.

\@ifpackagelater{xparse}{2018/04/17}{
\DeclareDocumentCommand \@printtoken { t^ ta m !ts } {%
  \IfBooleanTF {#1} 
               {\IfBooleanTF {#2}{\printtoken{A}{#3}~% 
                \IfBooleanTF {#4}
                 {\printtoken{p}{#3}}
                 {\printtoken{s}{#3}}}{%
                \IfBooleanTF {#4}
                 {\printtoken{P}{#3}}
                 {\printtoken{S}{#3}}}}{%
               {\IfBooleanTF {#2}{\printtoken{a}{#3}~}{}%
                \IfBooleanTF {#4}
                 {\printtoken{p}{#3}}
                 {\printtoken{s}{#3}}}}}
}{
\DeclareDocumentCommand \@printtoken { t^ ta m ts } {%
  \IfBooleanTF {#1} 
               {\IfBooleanTF {#2}{\printtoken{A}{#3}~% 
                \IfBooleanTF {#4}
                 {\printtoken{p}{#3}}
                 {\printtoken{s}{#3}}}{%
                \IfBooleanTF {#4}
                 {\printtoken{P}{#3}}
                 {\printtoken{S}{#3}}}}{%
               {\IfBooleanTF {#2}{\printtoken{a}{#3}~}{}%
                \IfBooleanTF {#4}
                 {\printtoken{p}{#3}}
                 {\printtoken{s}{#3}}}}}
}

% `\printtoken{switch}{token}` - print switch/token if defined within
% a `\tokenformat{...}`, produce warning if undefined defined, expand
% to `\tokenformat{\TEXTTOK@switch@token}`

\def\printtoken#1#2{%
  \ifcsname TOK@#1@#2\endcsname
    \tokenformat{\usetoken{#1}{#2}}\else
    \PackageWarning{open-logic-tokenize}{Token #2 undefined}{\color{red}#2}\fi}

% `\tokenformat` - wrapper to format token text; command can be
% redefined to eg print tokens in color. Warning: to work in headers
% it must be compatible with `\uppercase`

\def\tokenformat#1{#1}

% `\article{token}`, `\Article{token}` - just print the article for a
% token

\def\article#1{\printtoken{a}{#1}}
\def\Article#1{\printtoken{A}{#1}}

% use shorthand mechanism from `babel.sty` to define `!!` as shorthand
% for `\printtoken`

\useshorthands{!}
\defineshorthand{!!}{\@printtoken}

% OpenLogic Project
% Package for selective typesetting

% cribbed from tagging.sty package by Brent Longborough
% http://www.ctan.org/tex-archive/macros/latex/contrib/tagging 

% uses etoolbox's toggles for tags instead, allows for Boolean
% expressions in addition to tag lists

\NeedsTeXFormat{LaTeX2e}
\ProvidesFile{open-logic-selective.sty}

\RequirePackage{etoolbox}
\RequirePackage{xpunctuate} % for xcomma, xperiod

% initialize a list of tags as true

\newcommand{\tagtrue}[1]{%
  \def\do##1{\providetoggle{##1}\toggletrue{##1}}
    \docsvlist{#1}%
  \def\do##1{\providetoggle{not##1}\togglefalse{not##1}}
    \docsvlist{#1}%
}

% initialize a list of tags as false

\newcommand{\tagfalse}[1]{%
  \def\do##1{\providetoggle{##1}\togglefalse{##1}}%
    \docsvlist{#1}%
\def\do##1{\providetoggle{not##1}\toggletrue{not##1}}
    \docsvlist{#1}}

% tagcount -- count number of active tags in a tag list

\newcounter{tagnum}

\newcommand{\tagcount}[1]{%
  \setcounter{tagnum}{0}
  \def\do##1{\iftoggle{##1}{\stepcounter{tagnum}}{}}%
    \docsvlist{#1}}

% \iftag{tags}{text if tagged}{text if not tagged}

\newcommand{\iftagtext}{}
\newcommand{\iftag}[3]
  {\def\iftagtext{#3}%
   \def\do##1{%
     \iftoggle{##1}{\def\iftagtext{#2}}{}}%
   \docsvlist{#1}%
   \iftagtext}

% \tagrefs - output a list depending on tags.
\def\startcsvcomma{\def\csvcomma{}}
\def\tagpair#1/#2{\iftoggle{#1}{\edef\taglist@result{\taglist@result\csvcomma
      #2}\def\csvcomma{,,}}{}}
  
\newcommand{\tagrefs}[1]{\startcsvcomma%
  \def\taglist@input{#1}%
  \def\taglist@result{}%
  \def\do##1{\tagpair##1}%
  \docsvlist{#1}\cref{\taglist@result}}

% tagblock{tags} -- include contents of environment if a tag is true

\newenvironment{tagblock}[1]%
  {\iftag{#1}{%
      \let\comment\relax%
      \let\endcomment\relax%
    }{}%
    \comment\ignorespaces}% 
    {\endcomment\ignorespacesafterend}%

% \supressenv -- turn an envireonment into a comment 
% \excludeenv -- make the specified environment into a comment
% environment for the time being
% \includeenv -- restore a previously excluded environment

\def\supressenv#1{\begingroup\@bsphack
             \let\do\@makeother\dospecials\catcode`\^^M\active
             \let\verbatim@startline\relax
             \let\verbatim@addtoline\@gobble
             \let\verbatim@processline\relax
             \let\verbatim@finish\relax
             \def\@currenvir{#1}%
             \verbatim@}

\def\excludeenv#1{%
  \expandafter\let\csname orig@#1\expandafter
  \endcsname\csname #1\endcsname
  \expandafter\let\csname orig@end#1\expandafter
  \endcsname\csname end#1\endcsname
  \expandafter\def\csname #1\endcsname{\supressenv{#1}}
  \expandafter\let\csname end#1\endcsname\endgroup}
\def\includeenv#1{%
  \expandafter\let\csname #1\expandafter
  \endcsname\csname orig@#1\endcsname
  \expandafter\let\csname end#1\expandafter
  \endcsname\csname orig@end#1\endcsname}

% \tagitem<tags><tagged text><untagged text>: typeset an \item if tags active

\newcommand{\tagitem}[3]{\iftag{#1}{%
    \expandafter\ifstrempty{#2}{}{\item #2}}{
    \expandafter\ifstrempty{#3}{}{\item #3}}}

% \begin{tagenumerate}{tags} -- like enumerate, but only if more than
% one tag is active. If only one tag is active, just print the item
% without \item

\newenvironment{tagenumerate}[1]{%
  \tagcount{#1}%
  \ifnumgreater{\value{tagnum}}{1}{%
    \begin{enumerate}}{%
    \let\orig@tagitem\tagitem
    \let\orig@item\item
    \let\tagitem\iftag
    \let\item\relax}}{%
  \ifnumgreater{\value{tagnum}}{1}{%
    \end{enumerate}}{%
    \let\tagitem\orig@tagitem
    \let\item\orig@item}}

\def\startycommalist{\def\ycomma{\def\ycomma{, }}}

% tagprob environent -- include a problem, but only if tags in
% argument are true

\def\saveprob{\global\let\oldprob\prob
  \global\let\oldendprob\endprob}
\def\killprob{\global\let\prob\comment
  \global\let\endprob\endcomment}
\def\restoreprob{\global\let\prob\oldprob
  \global\let\endprob\oldendprob}

\newcommand{\tagprob}[2][tagTrue]{\saveprob\iftag{#1}{\iftag{#2}{}{\killprob}}{\killprob}}
\newcommand{\tagendprob}{\restoreprob}

%%
%% This is file `bussproofs-extra.sty',
%% generated with the docstrip utility.
%%
%% The original source files were:
%%
%% bussproofs-extra.dtx  (with options: `package')
%% 
%% This is a generated file.
%% 
%% Copyright (C) 2019 by Richard Zach <rzach@ucalgary.ca>
%% 
%% This file may be distributed and/or modified under the conditions of
%% the LaTeX Project Public License, either version 1.3 of this license
%% or (at your option) any later version.  The latest version of this
%% license is in:
%% 
%%    http://www.latex-project.org/lppl.txt
%% 
%% and version 1.3 or later is part of all distributions of LaTeX version
%% 1999/12/01 or later.
%% 
\NeedsTeXFormat{LaTeX2e}
\ProvidesPackage{bussproofs-extra}
   [2019/05/31 0.4 Extra commands for bussproofs.sty]

\RequirePackage{bussproofs}
\RequirePackage{tikz}
\newdimen\CenterCorrection
\newdimen\DiagCorrection
\newbox\myBoxLLL
\newbox\myBoxRLL
\tikzset{
    deduceLine/.style = {line width=1.1pt, loosely dotted}}

\def\straightDeduce{%
  \gdef\fDeduce{\tikz\draw[deduceLine] (0,0) -- (0,1);}
  \global\DiagCorrection=0pt
  \ignorespaces
}

\def\shortDeduce{%
  \gdef\fDeduce{\tikz\draw[deduceLine] (0,0) -- (0,.5);}
  \global\DiagCorrection=0pt
  \ignorespaces
}

\def\branchDeduce{%
  \gdef\fDeduce{\begin{tikzpicture}
      \draw[deduceLine] (0,0) -- (0,1);
      \draw[deduceLine] (-.5,.5) -- (0,0);
      \draw[deduceLine] (.5,.5) -- (0,0);
  \end{tikzpicture}}
  \global\DiagCorrection=0pt
  \ignorespaces
}

\def\ddotsDeduce{%
  \gdef\fDeduce{\begin{tikzpicture}
      \draw[deduceLine] (0,1) -- (1,0);
  \end{tikzpicture}}
  \setbox\myBoxA=\hbox{\fDeduce}
  \global\DiagCorrection=-\wd\myBoxA
  \ignorespaces
}

\def\dotsdDeduce{%
  \gdef\fDeduce{\begin{tikzpicture}
      \draw[deduceLine] (1,1) -- (0,0);
  \end{tikzpicture}}
  \setbox\myBoxA=\hbox{\fDeduce}
  \global\DiagCorrection=\wd\myBoxA
  \ignorespaces
}
\def\alwaysDeduce{\straightDeduce}
\straightDeduce
\def\Deduce$#1\fCenter#2${%
    \prepUnary%
    \buildConclusion{#1}{#2}%
    \setbox\myBoxA=\hbox{\fCenter}%
    \global\CenterCorrection=-.5\wd\myBoxA
    \joinDeduce%
    \resetInferenceDefaults%
    \ignorespaces%
}

\def\DeduceC#1{
    \prepUnary%
    \buildConclusionC{#1}%
    \global\CenterCorrection=0pt
    \joinDeduce%
    \resetInferenceDefaults%
    \ignorespaces%
}

\def\joinDeduce{%
    \global\advance\curCenter by -\hypKernAmt%
    \ifnum\curCenter<\newCenter%
        \displace=\newCenter%
        \advance \displace by -\curCenter%
        \kernUpperBox%
    \else%
        \displace=\curCenter%
        \advance \displace by -\newCenter%
        \kernLowerBox%
    \fi%
    \ifnum\DiagCorrection<0%
        \displace=-\DiagCorrection
        \kernLowerBox%
    \else
        \displace=\DiagCorrection
        \kernUpperBox%
    \fi%
    \advance\curCenter by-.5\DiagCorrection
    \buildDeduce%
    \buildScoreLabels%
    \ifx\rootAtBottomFlag\myTrue%
        \buildRootBottom%
    \else%
        \buildRootTop%
    \fi%
    \global \curScoreStart=\newScoreStart%
    \global \curScoreEnd=\newScoreEnd%
    \global \curCenter=\newCenter%
}

\def\buildDeduce{%
    \global\setbox \myBoxD =%
        \hbox{\fDeduce}%
    \displace = \wd\myBoxD % find width of vdots
    \global\curScoreStart = \curCenter%
    \global\advance\curScoreStart by -.5\displace%
    \global\curScoreEnd = \curCenter%
    \global\advance\curScoreEnd by .5\displace%
    \global\advance\curScoreStart by\CenterCorrection
    \global\advance\curScoreEnd by\CenterCorrection
}

\def\LeftLineLabel#1{%
  \global\def\displayLeftLineLabel{%
    {#1\hskip\labelSpacing}}
  \ignorespaces}

\def\RightLineLabel#1{%
  \global\def\displayRightLineLabel{%
    {\hskip\labelSpacing #1}}
  \ignorespaces}

\global\let\displayLeftLineLabel\relax
\global\let\displayRightLineLabel\relax

\def\LeftSubproofLabel#1{%
  \global\setbox\curBox =
  \hbox{\vbox to \ht\curBox{%
      \vfil
      \llap{#1$\left\{\vrule height .5\ht\curBox width 0pt\right.$}%
      \vfil}\box\curBox}%
}

\def\RightSubproofLabel#1{%
  \displace=\ht\curBox
  \global\setbox\curBox =
  \hbox{\box\curBox\vbox to \displace{%
      \vfil
      \rlap{$\left.\vrule height .5\displace width 0pt\right\}$#1}%
      \vfil}}%
}
\def\resetInferenceDefaults{%
    \global\def\theHypSeparation{\defaultHypSeparation}%
    \global\setbox\myBoxLL=\hbox{\defaultLeftLabel}%
    \global\setbox\myBoxRL=\hbox{\defaultRightLabel}%
    \global\def\buildScore{\alwaysBuildScore}%
    \global\def\theScoreFiller{\alwaysScoreFiller}%
    % reset line labels to nothing %bpextra
    \global\let\displayLeftLineLabel\relax %bpextra
    \global\let\displayRightLineLabel\relax %bpextra
    % reset to defaul deduce style %bpextra
    \alwaysDeduce %bpextra
    \gdef\hypKernAmt{0pt}% Restore to zero kerning.
}

\def\Axiom$#1\fCenter#2${%
    % Get level and correct names set.
    \prepAxiom%
    % Define the boxes
    % bpextra -- add line labels
    \setbox\myBoxA=\hbox{$\mathord{#1}\fCenter\mathord{\relax}$}%
    \setbox\myBoxB=\hbox{$#2$}% %bpextra
    \setbox\myBoxLLL=\hbox{\displayLeftLineLabel}% %bpextra
    \setbox\myBoxRLL=\hbox{\displayRightLineLabel}% %bpextra
    \global\setbox\curBox=%
         \hbox{\unhcopy\myBoxLLL%bpextra
         \hskip\ScoreOverhangLeft\relax
         \unhcopy\myBoxA
         \unhcopy\myBoxB
         \hskip\ScoreOverhangRight
         \unhcopy\myBoxRLL}%bpextra
    % Set the relevant dimensions for the boxes
    \global\curScoreStart=0pt \relax
    \global\curScoreEnd=\wd\curBox \relax
    \global\curCenter=\wd\myBoxA \relax %bpextra
    \global\advance \curCenter by \ScoreOverhangLeft%
    % bpextra adjust by dimensions of labels
    \global\advance \curCenter by \wd\myBoxLLL%bpextra
    \global\advance\curScoreStart by \wd\myBoxLLL%bpextra
    \global\advance\curScoreEnd by -\wd\myBoxRLL%bpextra
    % reset line labels to nothing %bpextra
    \global\let\displayLeftLineLabel\relax %bpextra
    \global\let\displayRightLineLabel\relax %bpextra
    \ignorespaces
}

\def\AxiomC#1{      % Note argument not in math mode
    % Get level and correct names set.
    \prepAxiom%
        % Define the box.
    \setbox\myBoxA=\hbox{#1}%
    \setbox\myBoxLLL=\hbox{\displayLeftLineLabel}% %bpextra
    \setbox\myBoxRLL=\hbox{\displayRightLineLabel}% %bpextra
    \global\setbox\curBox =%
        \hbox{\unhcopy\myBoxLLL%bpextra
            \hskip\ScoreOverhangLeft\relax%
            \unhcopy\myBoxA
            \hskip\ScoreOverhangRight\relax
            \unhcopy\myBoxRLL}% %bpextra
    % Set the relevant dimensions for the boxes
        \global\curScoreStart=0pt \relax
        \global\curScoreEnd=\wd\curBox \relax
        \global\curCenter=.5\wd\myBoxA \relax %bpextra
        \global\advance \curCenter by \ScoreOverhangLeft%
    % bpextra adjust by dimensions of labels
    \global\advance \curCenter by \wd\myBoxLLL%bpextra
    \global\advance\curScoreStart by \wd\myBoxLLL%bpextra
    \global\advance\curScoreEnd by -\wd\myBoxRLL%bpextra
    % reset line labels to nothing %bpextra
    \global\let\displayLeftLineLabel\relax %bpextra
    \global\let\displayRightLineLabel\relax %bpextra
    \ignorespaces
}

\def\buildConclusion#1#2{% Build lower sequent w/ center at \fCenter position.
    % Define the boxes
        \setbox\myBoxA=\hbox{$\mathord{#1}\fCenter\mathord{\relax}$}%
        \setbox\myBoxB=\hbox{$#2$}%
        \setbox\myBoxLLL=\hbox{\displayLeftLineLabel}% %bpextra
        \setbox\myBoxRLL=\hbox{\displayRightLineLabel}% %bpextra
    % Put them together in \myBoxC
    \setbox\myBoxC =%
          \hbox{\unhcopy\myBoxLLL%bpextra
          \hskip\ScoreOverhangLeft\relax%
        \unhcopy\myBoxA\unhcopy\myBoxB
        \hskip\ScoreOverhangRight
        \unhcopy\myBoxRLL}% %bpextra
    % Calculate the center of the \myBoxC string.
    \newScoreStart=0pt \relax%
    \newCenter=\wd\myBoxA \relax%
    \advance \newCenter by \ScoreOverhangLeft%
    \newScoreEnd=\wd\myBoxC%
    % bpextra adjust by dimensions of labels
    \global\advance\newCenter by \wd\myBoxLLL%bpextra
    \global\advance\newScoreStart by \wd\myBoxLLL%bpextra
    \global\advance\newScoreEnd by -\wd\myBoxRLL%bpextra
    % reset line labels to nothing %bpextra
    \global\let\displayLeftLineLabel\relax %bpextra
    \global\let\displayRightLineLabel\relax %bpextra
}

\def\buildConclusionC#1{% Build lower sequent w/o \fCenter present.
        % Define the box.
    \setbox\myBoxA=\hbox{#1}%
    \setbox\myBoxLLL=\hbox{\displayLeftLineLabel}% %bpextra
    \setbox\myBoxRLL=\hbox{\displayRightLineLabel}% %bpextra
\setbox\myBoxC =%
        \hbox{\unhcopy\myBoxLLL%bpextra
            \hskip\ScoreOverhangLeft\relax%
            \unhcopy\myBoxA
            \hskip\ScoreOverhangRight
            \unhcopy\myBoxRLL}%bpextra
    % Calculate kerning to line up centers
    \newScoreStart=0pt \relax%
    \newCenter=.5\wd\myBoxA \relax% bpextra
    \newScoreEnd=\wd\myBoxC%
        \advance \newCenter by \ScoreOverhangLeft%
    % bpextra adjust by dimensions of labels
    \global\advance\newCenter by \wd\myBoxLLL%bpextra
    \global\advance\newScoreStart by \wd\myBoxLLL%bpextra
    \global\advance\newScoreEnd by -\wd\myBoxRLL%bpextra
    % reset line labels to nothing %bpextra
    \global\let\displayLeftLineLabel\relax %bpextra
    \global\let\displayRightLineLabel\relax %bpextra
}
\endinput
%%
%% End of file `bussproofs-extra.sty'.
 % load version supplied by OLP
\ProvidesPackage{ptolemaicastronomy}

\tikzset{
  sphere/.style = {dotted},
  sphere intersection/.style = {fill=lightgray},
  sphere layer/.style = {fill=lightgray},
  proposition/.style={smooth,tension=1.7},
}

\pgfkeyssetvalue{/tikz/layerwidth}{.5}
\pgfkeyssetvalue{/tikz/innerfactor}{.4}

% \sphereplot{n} gives the plot codes for the nth sphere

\newcommand{\sphereplot}[1]{
  circle
    [radius=(#1)*\pgfkeysvalueof{/tikz/layerwidth}-
      \pgfkeysvalueof{/tikz/layerwidth}*\pgfkeysvalueof{/tikz/innerfactor}]
}


% \spheresystem[options]{n} draws a sphere system centered at
% the origin with n number of layers

\newcommand{\spheresystem}[2][]{
  \foreach \i in {1,...,#2}{
    \draw[sphere,#1] \sphereplot{\i} ;
  }
}

% \spherelayer[options]{n} shades the n-th layer

\newcommand{\spherelayer}[2][]{
  \begin{scope}[even odd rule]
    \fill[#1,sphere layer]
    \sphereplot{#2-1} \sphereplot{#2} ;
  \end{scope}
}

% \spherefill[options]{n} fills the nth sphere

\newcommand{\spherefill}[2][]{
    \fill[sphere intersection,#1]
    \sphereplot{#2} ;
}

% \sphereintersect[options]{path}{n} draws the path and shades the area
% of the path in the n th layer. Options only apply to the sphere layer

\newcommand{\sphereintersect}[3][]{
  \begin{scope}[even odd rule]
    \path[clip] #3;
    \spherefill[#1]{#2}
  \end{scope}
  \draw #3;
}

% \sphereproposition[options]{direction}{n}{width}{length} draws a
% proposition intersecting the n-th layer in direction away from the
% center of the sphere system, with endpoints length away from the
% vertex at an angle of width/2. Note that options applies to \draw,
% not to plot. To change the tension on the plot, e.g., say
% \tikzset{proposition/.style={smooth,tension=2}} -- in a group to
% make it apply only to a single proposition.

\newcommand{\propositionplot}[4]{
  plot [proposition]
  coordinates {+(#1+#3/2:#4)
    +(#1:#2*\pgfkeysvalueof{/tikz/layerwidth}-
    \pgfkeysvalueof{/tikz/layerwidth}*.9
    -\pgfkeysvalueof{/tikz/layerwidth}*\pgfkeysvalueof{/tikz/innerfactor})
    +(#1-#3/2:#4)}
}

\newcommand{\proposition}[5][]{
  \draw[proposition,#1] \propositionplot {#2}{#3}{#4}{#5} ;
  }

% \spherepropositionintersect does the same as \sphereproposition but
% also shades the are of intersection with the closest sphere

\newcommand{\propositionintersect}[5][]{
  \begin{scope}
  \path[clip] \propositionplot{#2}{#3}{#4}{#5};
  \spherefill[#1]{#3}
  \end{scope}
  \draw[proposition,#1] \propositionplot{#2}{#3}{#4}{#5};
}

% \spherepos[options]{direction}{n}{code} shifts the scope to a
% position in the center of the nth layer in direction angle from the
% center -- and then puts code there.

\newcommand{\spherepos}[4][]{
  \begin{scope}[#1,shift=(#2:#3*\pgfkeysvalueof{/tikz/layerwidth}-
      \pgfkeysvalueof{/tikz/layerwidth}/2-
        \pgfkeysvalueof{/tikz/layerwidth}*\pgfkeysvalueof{/tikz/innerfactor})]
    \path #4 ;
  \end{scope}
}


% Hooks
\def\OLEndChapterHook{}
\def\OLEndPartHook{}

% If compiled standalone as a subfile, execute \OLSubfileHook at beginning
\ifx\document\comment
\AtBeginDocument{\OLSubfileHook}
\fi
\let\OLSubfileHook\relax

% DeclareDocumentMacro -- for macros (no argument), doesn't check if defined

\newcommand*{\DeclareDocumentMacro}[2]{\def#1{#2}}

% We loading \oljobname-config.sty and \oljobname-envs.sty. Initialize
% \oljobname to \jobname -- this way you can use the same config files
% for differnt jobs.

\ifdefined\oljobname\relax\else\let\oljobname\jobname\fi

% Load configuration file, then custom changes if available

% % The Default OLP Configuration File `open-logic-config.sty`
% % OpenLogic Project
%
% Description
% ===========
%
% This file contains all commands and environments that are meant to
% be configured, changed, or adapted by a user generating their own
% text based on OLP text.  **Do not edit this file to customize your
% OLP-derived text!** A file `myversion.tex` adapted from
% `open-logic-complete.tex` (or from any of the contributed example
% master files) will include `myversion-config.sty` if it exists. It
% will do so after it loads this file, so your `myversion-config.sty`
% will redefine the defaults.  This means you won't have to include
% everythng, e.g., you can just change some tags and nothing else. You
% may copy and paste edfinitions you want to change into that file, or
% copy thi file, rename it `myversion-config.sty` and delete anything
% you'd like to leave as the default.

\NeedsTeXFormat{LaTeX2e}

% Symbols
% =======

% Formula metavariabes 
% --------------------
%
% Use the exclamation point symbol `!` immediately in front of an
% uppercase letter in math mode for formula metavariables.  By
% default, `!A`, `!B`, ... are typeset as $\varphi$, $\psi$, $\chi$,
% ... if you use the command `\olgreekformulas`. If this is not desired,
% and you'd like $A$, $B$, $C$, ... instead, use `\ollatinformulas`.

\olgreekformulas

% Logical symbols
% ---------------

% The following commands are used in the OLP texts for logical
% symbols. Their definitions can be customized to produce different
% output.

% ### Truth Values
%
% - `\True` defaults to $1$ and `\False` to $0$.

\DeclareDocumentMacro \True {1}
\DeclareDocumentMacro \False {0}

% ### Propositional Constants and Connectives
%
% - Falsity is `\lfalse` and defaults to $\bot$.

\DeclareDocumentMacro \lfalse {\bot}

% - Truth is `\ltrue` and defaults to $\top$.

\DeclareDocumentMacro \ltrue {\top}

% - Negation is `\lnot` and defaults to $\lnot$.  To use a different
% symbol (e.g., tilde), use the following line.

% `\DeclareDocumentMacro \lnot {\mathord{\sim}}`

% - Conjunction is `\land` and deaults to $\land$.  to use ampersand,
% uncomment the following line

% `\DeclareDocumentMacro \land {\mathbin{\&}}`

% - Disjunction `\lor` and defaults to $\lor$.

% - Conditional is `\lif` and defaults to $\rightarrow$.  To use a
% different symbol, replace `\rightarrow` in the definition, e.g., by
% `\supset`

\DeclareDocumentMacro \lif {\rightarrow}

% - The biconditional is `\liff` and defaults to $\leftrightarrow$.  To
% use the triple bar $\equiv$ replace with `\equiv`.

\DeclareDocumentMacro \liff {\leftrightarrow}

% Quantifiers 
% ----------- 

% The quantifier symbols are provided as commands `\lexists` and
% `\lforall` which take two optional arguments. If no arguments are
% provided, it they just typeset the quantifier symbol. With one
% optional argument they produce the quantifier together with a
% variable, and this may include parenthesesaround the quantifier and
% variable. The second optional argument producesthe
% quantifier/variable combination plus the formula in the scope of the
% formula with appropriate spacing.  For instance,
% `\lexists[x][!A(x)]` will, bydefault, produce $\exists
% x\,\varphi(x)$.

% - The existential quantifier is `\lexists`.  Replace `\exists` with
% `\boldsymbol{\exists}` for boldface, or redefine appropriately if
% you want parentheses around $\exists x$.

\DeclareDocumentCommand \lexists { o o } {
    \IfNoValueTF {#1} 
        { \exists }    % no arguments
        { \exists #1 } % one argument: variable
    \IfNoValueTF {#2} 
        \relax
        { \, #2 }      % two arguments: space and matrix
}

% - The universal quantifier is `\lforall`.

\DeclareDocumentCommand \lforall { o o } {
    \IfNoValueTF {#1} 
        { \forall }    % no arguments
        { \forall #1 } % one argument: variable
    \IfNoValueTF {#2} 
        \relax
        { \, #2 }      % two arguments: space and matrix
}

% - The identity relation is also provided as `\eq`. By itself, it
% produces the identity reation symbol (default: $=$) by itself. With
% two optional arguments, it typesets the corresponding atomic
% formula, e.g., `\eq[x][y]` produces $x = y$.  `\eq/` produces the
% negated symbol (formula).

\DeclareDocumentCommand \eq { t{/} o o } {
  \IfNoValueTF {#3}
    % no optional arguments: just typeset symbol
    { \IfBooleanTF{#1}{ \neq }{ = } }
    % optional arguments: typeset atomic formula
    { \IfBooleanTF{#1}{ #2 \neq #3}{#2 = #3} }
}

% - The sequent symbol `\Sequent` produces $\Rightarrow$ by
% default. Change the definition for $\vdash$, or another symbol.

\DeclareDocumentMacro \Sequent {\Rightarrow}

% The sequent symbol in proofs displays as the above sequent symbol

\DeclareDocumentMacro \fCenter {\ensuremath{\,\Sequent\,}}


% Metalogical Relations
% ---------------------
%
% Metalogical relationships, such as truth in a structure, validity,
% consequence, and provability, are also provided as commands. Uniform
% use of these commandsinstead of hard-coded typesetting according to
% specific conventions guarantees that by changing the definitions
% below you can uniformly change notation in the text.


% ### Substitution

% `Subst{!A}{t}{x}`: The operation of substituting a term for a (free)
% variable in another term or in a formula.  Common notations are
% $\varphi^t_x$, $\varphi(t/x)$, $\varphi[x/t]$, or $S^t_x \varphi$.

\DeclareDocumentCommand \Subst { m m m } {
  #1[#2/#3]}



% ### The satisfaction/truth relation 

% `\Sat[/]{M}{!A}[s]` is the relation of being satisfied in a
% structure (relative to an assignment) isprovided as the command
% `\Sat` with two mandatory arguents (the structure and the formula)
% and one optional argument (the assignment).  Use `\Sat/` to create
% the negated relation.  By default, `\Sat{M}{!A}[s]` is typeset as
% $\mathfrak{M}, s \models \varphi$.

\DeclareDocumentCommand \Sat { t{/} m m o } {
  \IfBooleanTF{#1}{
    % negated
    \IfNoValueTF {#4} 
        { \Struct #2 \not\models #3 }
        { \Struct #2, #4 \not\models #3}}{
    % not negated
    \IfNoValueTF {#4} 
        { \Struct #2 \models #3 }
        { \Struct #2, #4 \models #3 }}
}


% ### The derivability relation

% `\Proves` is used to create the symbol for the derivability
% relation, `\Proves/` for the negation. By default this creates
% $\vdash$. The optional argument is used for the calculusor logic
% relative to which the provability relation is defined; by default it
% creates a subscript on the turnstile.

\DeclareDocumentCommand \Proves { t{/} o } {
  \IfBooleanTF {#1}{
    \IfNoValueTF {#2} 
        { \not\vdash }
        { \not\vdash_{#2} }}{
    \IfNoValueTF {#2} 
        { \vdash }
        { \vdash_{#2} }}
}

% ### The semantic consequence relation relation

% `\Entails` is the semantic counterpart of `\Proves`. 

\DeclareDocumentCommand \Entails { t{/} o } {
  \IfBooleanTF {#1}{
    \IfNoValueTF {#2} 
        { \not\models }
        { \not\models_{#2} }}{
    \IfNoValueTF {#2} 
        { \models }
        { \models_{#2} }}
}


% ### Model-theoretic notions and symbols

% - `\Domain{M}` - domain of a structure, e.g., `\Domain M` gives
% $\left|\mathfrak M\right|$.

\DeclareDocumentCommand \Domain { m }{\left| \Struct #1 \right|}

% - `\Assign{R}{M}` - Assignment (value of) of a constant/predicate symbol
% in a structure; e.g., `\Assign R M` produces $R^\mathfrak{M}$.

\DeclareDocumentCommand \Assign { m m }{#1^{\Struct #2}}

% - `\Value{t}{M}[s]` - Value of a term in a structure. Takes two mandatory
% arguments (term and structure) and one optional argument (variable
% assignment). By default, `\Value{t}{M}[s]` produces 
% $\mathrm{Val}^\mathfrak{M}_s(t)$.

\DeclareDocumentCommand \Value { m m o} {
    \IfNoValueTF {#3}
        % optional argument not present
        { \mathrm{Val}^{\Struct #2}(#1) }
        % optional argument present
        { \mathrm{Val}^{\Struct #2}_{#3}(#1) }
}

% - `\substruct`: symbol for the substructure relation

\DeclareDocumentMacro \substruct {\subseteq}

% - `\Theory{M}`: theory of a structure

\DeclareDocumentCommand \Theory { m } {\mathrm{Th}(\Struct{#1})}

% - `\Mod[L](L'){T}`: class of models of a theory/sentence $T$ in a
% language $L$ and logic $L'$.

\DeclareDocumentCommand \Mod { o d() m } {
  \IfNoValueTF {#2} {
  % optional logic argument not present
    \IfNoValueTF {#1}{
      \mathrm{Mod}(#3) }{
      \mathrm{Mod}^{\Lang{#1}}(#3) }}{
  % optional logic argument present
    \IfNoValueTF {#1}{
      \mathrm{Mod}_{#2}(#3)}{
      \mathrm{Mod}_{#2}^{\Lang{#1}}(#3)}}
}


% - `\elemequiv`: elementary equivalence (infix relation)

\DeclareDocumentCommand \elemequiv { t{/} o } {
  \IfBooleanTF {#1}{
    \IfNoValueTF {#2} 
        { \not\equiv }
        { \not\equiv_{#2} }}{
    \IfNoValueTF {#2} 
        { \equiv }
        { \equiv_{#2} }}
}

% - `\iso[/]`: relation of (partial) isomorphism

\DeclareDocumentCommand \iso { t{/} o } {
  \IfBooleanTF {#1}{
    \IfNoValueTF {#2} 
        { \not\simeq }
        { \not\simeq_{#2} }}{
    \IfNoValueTF {#2} 
        { \simeq }
        { \simeq_{#2} }}
}



% - `\ident`: syntactic identity between expressions (infix relation),

\DeclareDocumentMacro \ident {\equiv}

% - `\QuantRank`: quantifier rank of a formula

\DeclareDocumentCommand \QuantRank { m } {\mathrm{qr}(#1)}

% `\Expan{M}{R}`: expansion of a structure by a relation (etc.)

\DeclareDocumentCommand \Expan { m m } {(\Struct{#1}, #2)}


% Recursion-theoretic Notions and Symbols
% ---------------------------------------

% - `\Proj{n}{i}`: projection functions

\DeclareDocumentCommand \Proj { m m } {P^{#1}_{#2}}

% - `\tsub`: truncated subtraction function

\DeclareDocumentMacro \tsub {\mathbin{\dot-}}

% - `\Char{R}`: characteristic function

\DeclareDocumentCommand \Char { m } {\chi_{#1}}

% - `\defis`: definitional identity

\DeclareDocumentMacro \defis {=} %{\mathrel{=_\mathrm{df}}}

% - `\defiff`: definitional equivalence

\DeclareDocumentMacro \defiff {\Leftrightarrow}

% - `\concat`: concatenation of sequences

\DeclareDocumentMacro \concat {\frown}

% - `\umin{x}{!A}`: unbounded minimization

\DeclareDocumentCommand \umin { m m } {\mu #1 \; #2}

% - `\bmin{x < y}{!A}`: bounded minimization

\DeclareDocumentCommand \bmin { m m } {\fn{min} \; #1 \; #2}

% - `\cfind{e}[n]`: partial computable function with index $e$

\DeclareDocumentCommand \cfind { m o } {%
    \IfNoValueTF {#2}
        % optional argument not present
        { \varphi_{#1} }
        % optional argument present
        { \varphi_{#1}^{#2} }
}

% - `\redone`: one-step reduction

\DeclareDocumentMacro \redone {\mathrel{\triangleright_1}}

% - `\red`: reduction

\DeclareDocumentMacro \red {\triangleright}

% - `\lambd[x][!A]`: lambda abstract

\DeclareDocumentCommand \lambd { o o } {
    \IfNoValueTF {#1} 
        { \lambda }    % no arguments
        { \lambda #1 } % one argument
    \IfNoValueTF {#2} 
        \relax
        { .\, #2 }      % two arguments
}

% - `\num{n}` : numeral corresponding to a number

\DeclareDocumentCommand \num { m } {\overline{#1}}

% - `\scode{s}`: code for a symbol

\DeclareDocumentCommand \scode { m } {\fn{c}(#1)}

% - `\Gn{!A}`: G\"odel number of a string of symbols

\DeclareDocumentCommand \Gn { m } {\#(#1)}

% - `\mSat{M}{!A}[w]`: modal stasfaction relation

\DeclareDocumentCommand \mSat { t{/} m m o } {
  \IfBooleanTF{#1}{
    % negated
    \IfNoValueTF {#4} 
        { \Struct #2 \nVdash #3 }
        { \Struct #2, #4 \nVdash #3}}{
    % not negated
    \IfNoValueTF {#4} 
        { \Struct #2 \Vdash #3 }
        { \Struct #2, #4 \Vdash #3 }}
}


% Special Sets and Mathematical Symbols
% -------------------------------------

% ### Set-theoretic operators

% - Set abstracts: Use `\Setabs{x}{!A(x)}` to produce the set abstract
% $\{ x : \varphi(x) \}$. If you prefer a $\mid$ to :, change the
% definition accordingly.

\DeclareDocumentCommand \Setabs { m m }{\{ #1 : #2 \}}

% - `\Pow{X}`: Power set, produces $\wp(X)$

\DeclareDocumentCommand \Pow { m }{\wp(#1)}

% - `\dom{f}`: domain of a function

\DeclareDocumentCommand \dom { m }{\fn{dom}(#1)}

% - `\ran{f}`: range of a function
\DeclareDocumentCommand \ran { m }{\fn{ran}(#1)}

% - `\len{s}`: length of a sequence
\DeclareDocumentCommand \len { m }{\fn{len}(#1)}

% - `\emptyseq`: the empty sequence
\DeclareDocumentMacro \emptyseq {\emptyset}

% - `\restrict`: restriction of a function to a set (infix operator)
\DeclareDocumentMacro \restrict {\upharpoonright}

% - `\Complement{X}`: complement of a set
\DeclareDocumentCommand \Complement { m } {\overline{#1}}

% - `\card{X}`: cardinality of a set
\DeclareDocumentCommand \card { m } {\left| #1 \right|}

% - `\tuple{x,y}`: pairs, tuples, sequences
\DeclareDocumentCommand \tuple { m } {\langle #1 \rangle}

% - `\comp{f}{g}`: composition of f with g, defaults to $g \circ f$
\DeclareDocumentCommand \comp { m m }{g \circ f}

% - `\pto`: partial function arrow

\DeclareDocumentMacro \pto {\mathrel{\ooalign{\hfil$\mapstochar\mkern
      5mu$\hfil\cr$\to$}}}

% - `\defined`, `\undefined`: postfix for defined, undefined
% functions

\DeclareDocumentMacro \defined {\downarrow}
\DeclareDocumentMacro \undefined {\uparrow}

% ### Particular sets

% - Natural numbers: `\Nat`
\DeclareDocumentMacro \Nat {\mathbb{N}}

% - Integers: `\Int`
\DeclareDocumentMacro \Int {\mathbb{Z}}

% - Real numbers: `\Real`
\DeclareDocumentMacro \Real {\mathbb{R}}

% - Rational numbers: `\Rat`
\DeclareDocumentMacro \Rat {\mathbb{Q}}

% - The set $\{0, 1\}$: `\Bin`
\DeclareDocumentMacro \Bin {\mathbb{B}}

% - Identity relation: `\Id{X}`
\DeclareDocumentCommand \Id { m } {\mathord{\mathrm{Id}_{#1}}}

% ### Symbols for Turing Machines

% - `\TMendtape` - symbol indicating left end of tape
\DeclareDocumentMacro \TMendtape {\triangleright}

% - `\TMblank` - symbol for a blank
\DeclareDocumentMacro \TMblank {\mathord{\text{\textvisiblespace}}}

% - `\TMstroke` - single stroke symbol on tape
\DeclareDocumentMacro \TMstroke {\mid}

% - `\TMright` - symbol for move right instruction
\DeclareDocumentMacro \TMright {R}

% - `\TMleft` - symbol for move left instruction
\DeclareDocumentMacro \TMleft {L}

% - `\TMstay` - symbol for the stay instruction
\DeclareDocumentMacro \TMstay {N}

% ### Functions and Function/Relation symbols

% - `\Part`: the parthood predicate 

\DeclareDocumentCommand \Part { m m } {\Atom{\Obj P}{#1, #2}}

% - `\Prf`: the proof relation

\DeclareDocumentCommand \Prf { o } { \mathrm{Pr}\IfNoValueTF {#1} {} {_\Th{#1}}}
\DeclareDocumentCommand \OPrf { o } { \mathsf{Pr}\IfNoValueTF {#1} {} {_\Th{#1}}}

% - `\Prov`: the provability relation

\DeclareDocumentCommand \Prov { o } { \mathrm{Prov}\IfNoValueTF {#1} {} {_\Th{#1}}}
\DeclareDocumentCommand \OProv { o } { \mathsf{Prov}\IfNoValueTF {#1} {} {_\Th{#1}}}

% - `\RProv`: the Rosser provability relation

\DeclareDocumentCommand \RProv { o } { \mathrm{RProv}\IfNoValueTF {#1} {} {_\Th{#1}}}
\DeclareDocumentCommand \ORProv { o } { \mathsf{RProv}\IfNoValueTF {#1} {} {_\Th{#1}}}

% - `\OCon`: the consistency statement

\DeclareDocumentCommand \OCon { o } { \mathsf{Con}\IfNoValueTF {#1} {} {_\Th{#1}}}

% Typesetting commands for logical concepts
% =========================================

% In order to unformly typeset certain types of symbols uniformly, the
% OLP texts use special commands which carry out the
% typesetting. Thus, e.g., all structures, which appear in the texts,
% say, as `\Struct{M}`, can be typeset according to preference. By
% default, they are typeset in fraktur (e.g., $\mathfrak M $) but this
% can be configured by changing the definition of the `\Struct`
% command.  Note that often the default behavior is to only apply the
% typeface command to the first token in the argument, e.g., `\Struct
% M_n` will generate $\mathfrak M_n$ and not $\mathfrak{M_n}$.

% - `\Struct{M}` - First-order structures; by default, the first token
% in Fraktur

\DeclareDocumentCommand \Struct { m }{\mathfrak #1}

% - `\mStruct{M}` - modal structures; default: set first token in
% Fraktur

\DeclareDocumentCommand \mStruct { m }{\mathfrak #1}

% - `\Lang{L}` - Languages; default: set first token in callgraphic
% font

\DeclareDocumentCommand \Lang { m }{\mathcal #1}

% - `\Log{L}` - Logics; default: set entirely in boldface

\DeclareDocumentCommand \Log { m }{\mathbf{#1}}

% - `\Obj` - Object-language symbols; default: set entirely in
% sans-serif italics

\DeclareDocumentCommand \Obj { m }{\mathsfit{#1}}

%  - `\Atom{P}{t_1, t_2)` - Atomic formula or term; default produces
%  predicate symbol followed by argumnets surrounded by
%  parentheses. Some prefer no parentheses around the arguments.

\DeclareDocumentCommand \Atom { m m }{ \mathord{#1}(#2) }

% - `\PIso` - Set of all partial isomorphisms

\DeclareDocumentCommand \PIso { m }{\mathcal{#1}}

% `\fn{func}` -- typeset a function name

\DeclareDocumentCommand \fn { m } {\mathrm{#1}}

% `\Th{T}` -- typeset name of a theory

\DeclareDocumentCommand \Th { m } {\mathbf{#1}}

% Sets of Expressions
% ===================
%
% The following commands provide uniform notation for sets of
% expressions, such as the set of formulas of a language. Typically
% the parameter (language) can be providedas an optional argument.  By
% default, if the argument is given, they aretypeset following the
% symbol within parentheses, i.e., the symbols are treated as
% operators (e.g., $\mathrm{Frm}(\mathcal{L})$. Redefine the commands
% to, e.g., typeset the parameters as subscripts (e.g.,
% $\mathrm{Frm}_\mathcal{L}$.

% - `\Var`: the set of propositional variables

\def\Var{\mathrm{Var}}

% - `\Trm[L]`: the set of terms (of a language)

\DeclareDocumentCommand \Trm { o } {
    \IfNoValueTF {#1} 
        { \mathrm{Trm} }
        { \mathrm{Trm}({\Lang #1}) }
}

% - `\Frm[L]`: the set of formulas (of a language)
\DeclareDocumentCommand \Frm { o } {
    \IfNoValueTF {#1} 
        { \mathrm{Frm} }
        { \mathrm{Frm}({\Lang #1}) }
}

% - `\SubFrm{!A}`: the set of subformulas of a formula (mandatory
% argument: formula)

\DeclareDocumentCommand \SubFrm { m } {
        \mathrm{SFrm}({#1})
}

% - `\Sent[L]`: the set of sentences (of a language)
\DeclareDocumentCommand \Sent { o } {
    \IfNoValueTF {#1} 
        { \mathrm{Sent} }
        { \mathrm{Sent}({\Lang #1}) }
}

% Commands for uniform formulations
% =================================

% Inductive definitions
% ---------------------

% Inductive definitions typically divide into cases depening on the
% form of a formula for which a concept if defined.  To uniformly
% typeset these definitions flexibly, we provide a command
% `\indcase{formula}{complex formula}{case text}` which (a) typesets a
% uniform description of a case and (b) defines `\indfrm` and
% `\indcomplex` which can then be used in the definitions. For
% instance, in an inductive definition of $\models A$, you might say
% `\indcase{A}{B \land C}{$\models \indfrm$ iff $\models B$ and
% $\models C$}` to produce: "If $A \equiv B \land C$, then $\models A$
% iff $\models B$ and $\models C$" or, alternatively and more
% succinctly: "$\models B \land C$ iff $\models B$ and $\models C$".
% Use the starred version for the atomic case, and a ! instead of the
% star for a case you to leave as an exercise.

\DeclareDocumentCommand \indcase { s t{!} m m +m }{%
  \DeclareDocumentMacro \indfrm {#3}%
  \DeclareDocumentMacro \indfrmp {#3}%
  \DeclareDocumentMacro \indcomplex {#4}%
  \IfBooleanTF{#1}
     {$#3$ is atomic: }{$#3 \ident #4$: }
  \IfBooleanTF{#2}
     {exercise.}
     {#5}}

% Tokens
% ======

% The following terms are *tokenized* throughout OLP. This means that
% by changing the definition in the configuration file, the term as
% printed will also change. This is simpler than searching and
% replacing these terms in all OLP texts, and it will also treat
% plurals as well as occurrences of the term at the beginning of a
% sentence (where it should be capitalized) correctly.

% Tokens are defined using the `\settexttoken` command.
% ```
% \settexttoken{token}{singular}{plural}[Singular][Plural]
% ```
% Here `token` is the term as it is used in the source text, where it
% typically is used as `!!{token}`.  `singular` and `plural` are the
% text you want printed wherever an OLP text contains `!!{token}` or
% `!!{token}s`.  In sentence-initial position, OLP texts would use
% `!!^{token}` and `!!^{token}s` to create capitalized versions of
% these token replacements.  By default, they are generated from
% `singluar` and `plural` by capitalizing the first chracter, but can
% be provided explicitly to `\settexttoken` as optional arguments.
% `!!a{token}` produces an indefinite article plus the token
% replacement. This is usually "a" unless the settexttoken command is
% used with a star, as in `\settexttoken{element}*{element}{elements}`,
% in which case it produces "an". This can be combined with `^` to
% produce the uppercase version, e.g., `!!^a{element}` for "An
% element".  `\article{token}` and `\Article{token}` produce just the
% article by itself (lower or uppercase, respectively).

% - `language`: defaults to ``language'', redefine for
% ``signature''

\settexttoken{language}{language}{languages}

% - `formula`: defaults to "formula/formuls", redefine for
% plural "formulae," or for "wff".

\settexttoken{formula}{formula}{formulas}

% - `subformula`: defaults to "subformula", redefine for plural
% "subformulae," hyphenated spelling "sub-formula", or "sub-wff".

\settexttoken{subformula}{subformula}{subformulas}

% - `sentence`: defaults to "sentence", redefine for
% "closed formula" etc.

\settexttoken{sentence}{sentence}{sentences}

% - `variable`: defaults to "variable/variables", redefine to be more
% specific, e.g., "individual variable", "object veriable".

\settexttoken{variable}{variable}{variables}

% - `constant`: defaults to "constant", redefine for "individual
% constant'', "constant symbol."

\settexttoken{constant}{constant symbol}{constant symbols}

% - `predicate`: defaults to "predicate symbol".

\settexttoken{predicate}{predicate symbol}{predicate symbols}

% - `function`: defaults to "function symbol".

\settexttoken{function}{function symbol}{function symbols}

% - `operator`: defaults to "logical operator", redefine for "connective".

\settexttoken{operator}{logical operator}{logical operators}

% - `main operator`: defaults to "main operator", redefine for
% "outermost operator".

\settexttoken{main operator}{main operator}{main operators}

% - `free for`: defaults to "free for", redefine for
% "substitutable for" as in Enderton.

\settexttoken{free for}{free for}{free for}

% - `identity`: defaults to "identity predicate", redefine for
% "equality predicate."

\settexttoken{identity}*{identity predicate}{identity predicates}

% - `conditional`: defaults to "conditional", redefine for "implication."

\settexttoken{conditional}{conditional}{conditionals}

% - `biconditional`: defaults to "biconditional", redefine for "equivalence."

\settexttoken{biconditional}{biconditional}{biconditionals}

% - `falsity`: the falsity symbol, defaults to "falsity", redefine for
% "absurdity."

\settexttoken{falsity}{falsity}{falsities}

% - `truth`: the falsity symbol, defaults to "truth", redefine for
% "verum" or "top".

\settexttoken{truth}{truth}{truth}

% - `structure`: term for first-order structures, defaults to
% "structure", redefine for "interpretation", "model".

\settexttoken{structure}{structure}{structures}

% `domain`: domain of a structure

\settexttoken{domain}{domain}{domains}

% `\value`: value (denotation) of a term

\settexttoken{value}{value}{values}

% `derivation`: derivation in a calculus, proof

\settexttoken{derivation}{derivation}{derivations}

% `derive`: derive in a calculus, prove

\settexttoken{derive}{derive}{derives}

% `derivable`: derivable in a calculus, provable

\settexttoken{derivable}{derivable}{derivable}

% `derivability`: derivability in a calculus, provability

\settexttoken{derivability}{derivability}{derivabilities}

% `nonderivability`: derivability in a calculus, unprovability

\settexttoken{nonderivability}{non-derivability}{non-derivabilities}


% - `enumerable`: term for finite or countably infinite; defaults to
% "enumerable", redefine for "countable".

\settexttoken{enumerable}*{enumerable}{enumerable}

% - `nonenumerable`: term for uncountable; defaults to
% "non-enumerable", redefine for "uncountable".

\settexttoken{nonenumerable}{non-enumerable}{non-enumerable}

% - `denumerable`: term for countably infinite; defaults to "denumerable". 

\settexttoken{denumerable}{denumerable}{denumerable}

% - `element`: element of a set; redefine for "member"

\settexttoken{element}*{element}{elements}

% - `injective`, `injection`: redefine for "one-one"

\settexttoken{injective}*{injective}{injective}
\settexttoken{injection}*{injection}{injections}

% - `surjective`, `surjection`: redefine for "onto"

\settexttoken{surjective}{surjective}{surjective}
\settexttoken{surjection}{surjection}{surjections}

% - `bijective`, `bijection`: redefine for "one-one onto"

\settexttoken{bijective}{bijective}{bijective}
\settexttoken{bijection}{bijection}{bijections}

% Tags
% ====

% Tags are used to guide selective compilation, using
% `sty/open-logic-selective.sty`.  Tags are initialized using
% `\tagtrue` and `\tagfalse`, which also initialize, for each `tag`
% given as argument, a corresponding tag `nottag` with opposite truth
% value is automatically provided. 

% - Primitives: tags for (sets of) operators which are treated as
%   primitives. Default: all as primitives. 

\tagtrue{prvNot,prvOr,prvAnd,prvIf,prvIff,prvTrue,prvFalse,prvEx,prvAll}

% - Defined operators: tags for operators which are defined, not
%  primitive. Default: none.  Note: Not all combinations of
%  primitive/defined operators will result in complete definitions!

\tagfalse{defNot,defOr,defAnd,defIf,defIff,defTrue,defFalse,defEx,defAll}

% - Cases in proofs: tags for (sets of) operators for which cases are
%   proved; if the corresponding tag is off then the case is added as
%   a problem. Default: prove all cases. Note: propositional constants
%   are part of induction base case and are always treated, if
%   included.

\tagfalse{probNot,probOr,probAnd,probIf,probIff,probEx,probAll}

% - Limit clause in inductive definitions: do you want inductive
% definitions to have a "Nothing else is a ..." clause?

\tagtrue{limitClause}

% tagTrue - a true tag respectively.

\tagtrue{tagTrue}

% TMs - Turing machines have been discussed

\tagtrue{TMs}

% lambda - Lambda calculus has been discussed

\tagtrue{lambda}

\IfFileExists{\oljobname-config.sty}
             {\input{\oljobname-config.sty}}
             {}

% Load environment definition file, or custom file if available

\IfFileExists{\oljobname-envs.sty}
             {\input{\oljobname-envs.sty}}
             {% % The Default OLP Theorems File `open-logic-envs.sty`
% % OpenLogic Project
%
% Description
% ===========
%
% This file contains all theorem and other environments that are meant to
% be configured, changed, or adapted by a user generating their own
% text based on OLP text.  **Do not edit this file to customize your
% OLP-derived text!** A file `myversion.tex` adapted from
% `open-logic-complete.tex` (or from any of the contributed example
% master files) will include `myversion-thm.sty` if it exists, and
% otherwise will include this file. So if you'd like to make changes,
% instead **copy this file to `myversion-config.sty` and edit for
% customization**.

\NeedsTeXFormat{LaTeX2e}


% Theorem-like Environments
% -------------------------

% Theorem-like environments are provided for theorem (`thm`),
% definition (`defn`), lemma (`lem`), proposition (`prop`), corollary
% (`cor`), example (`ex`), problem (`prob`), remark (`rem`), and
% convention (`conv`). The definitions make use of the `thmtools`
% package.  To change the numbering or style of these environments,
% adjust their definitions. Refer to the [`thmtools`
% documentation](http://www.ctan.org/pkg/thmtools).

\declaretheorem[
  style=plain,
  name=Theorem,
  numberwithin=chapter]{thm}

\declaretheorem[
  style=definition,
  name=Example,
  sibling=thm]{ex}

\declaretheorem[
  style=plain,
  name=Lemma,
  refname={Lemma,Lemmata},
  sibling=thm]{lem}

\declaretheorem[
  style=plain,
  name=Proposition,
  sibling=thm]{prop}

\declaretheorem[
  style=plain,
  name=Corollary,
  refname={Corollary,Corollaries},
  sibling=thm]{cor}

\declaretheorem[
  style=definition,
  name=Definition,
  sibling=thm]{defn}

\declaretheorem[
  style=definition,
  name=Problem,
  numberwithin=chapter]{prob}

% A wrapper to allow problems with tags, use \begin{probtag}{tags}
\newenvironment{probtag}[1]{\begin{prob}}{\end{prob}}

\declaretheorem[
  style=remark,
  name=Remark]{rem}

\declaretheorem[
  style=remark,
  name=Note]{note}

\declaretheorem[
  style=remark,
  name=Case]{case}

\declaretheorem[
  style=remark,
  name=Convention]{conv}

% These are the names of the environments that `\cref` adds
% automatically. `thmtools` should do that as well, but sometimes it
% doesn't work, so we set them explicitly.  Redefine, e.g., if you
% want references to "Prop. 5.4"

\crefname{thm}{Theorem}{Theorems}
\crefname{ex}{Example}{Examples}
\crefname{defn}{Definition}{Definitions}
\crefname{lem}{Lemma}{Lemmata}
\crefname{prop}{Proposition}{Propositions}
\crefname{prob}{Problem}{Problems}
\crefname{rem}{Remark}{Remarks}
\crefname{figure}{Figure}{Figures}




% Other Enviroments
% -----------------

% OLP texts make use of a number of environments to encapsulate types
% of discursive text.  By default, these environments simply print
% their content without any special treatment. If you want to typeset
% any of them differntly, you can change the definition of the
% environment here. E.g., you might want digressions in asmaller font
% and indented.  Refer to the \LaTeX\ documentation for how to
% accomplish this.

% OLP will provide two custom mechanisms for treating environments:
% supressing and deferring. Environments can be supressed by
% redefining them as the `comment` environment. For instance, to not
% print any digressions at all, add the lines
% ```
% \let\digress\comment
% \let\enddigress\endcomment
% ```
% after the definition of `digress`.

% - `explain`: Any long and/or informal explanation that's useful for
% novice readers but could be left out

\newenvironment{explain}{}{}

% - `intro`: for comments and comparisions to other introductory
% texts, e.g., regarding terminology, conventions, or proof methods.

\newenvironment{intro}{}{}

% - `pedantic`: for pedantic comments which users may want to exclude.

\newenvironment{pedantic}{}{}

% - `digress`: for digressions

\newenvironment{digress}{}{}

% - `history`: for historical notes and remarks

\newenvironment{history}{}{}
}

% olfilepath -- holds the path of the currently processed file
% relative to the main input file.

% olimport -- import a text. First optional argument is subdirectory,
% if present subimport #1/#2, no optional argument: just \subfile it.
% Starred version interprets #1 as relative to main O-L directory
% \olpath/content/.  When asking for a file from a subdirectory, change
% \olfilepath accordingly. This happens within a gorup, so \olfilepath
% reverts to its original content after the \subfile.  Second optional
% argument is name of sectioning command to use for \olsection.  Use
% ``nosection'' to supress sectioning, ``subsection'' to make it a
% subsection instead of a section, etc..

\def\olfilebase{}
\edef\olfilename{\jobname.tex}

\DeclareDocumentCommand \olimport { s o m o }{
  \begingroup
  \IfNoValueTF{#4}{%
    \def\ol@sectioncs{section}}{%
    \def\ol@sectioncs{#4}}
  \IfNoValueTF{#2}{%
    \def\olfilename{#3.tex}%
    \subfile{\olfilebase #3}%
}{%
    \def\olfilename{#3.tex}%
    \IfBooleanTF{#1}{%
      \edef\olfilebase{\olpath content/#2/}}{%
      \edef\olfilebase{\olfilebase #2/}}%
    \subfile{\olfilebase #3}}
  \endgroup}

% Goedel corner quotes

\newbox\gnBoxA
\newdimen\gnCornerHgt
\setbox\gnBoxA=\hbox{$\ulcorner$}
\global\gnCornerHgt=\ht\gnBoxA
\newdimen\gnArgHgt
\def\gn #1{%
\setbox\gnBoxA=\hbox{$#1$}%
\gnArgHgt=\ht\gnBoxA%
\ifnum     \gnArgHgt<\gnCornerHgt \gnArgHgt=0pt%
\else \advance \gnArgHgt by -\gnCornerHgt%
\fi \raise\gnArgHgt\hbox{$\ulcorner$} \box\gnBoxA %
\raise\gnArgHgt\hbox{$\urcorner$}}

% Additional symbols

\DeclareSymbolFont{symbolsC}{U}{ntxsyc}{m}{n}
\DeclareMathSymbol{\boxright}{\mathbin}{symbolsC}{128}
\DeclareMathSymbol{\diamondright}{\mathbin}{symbolsC}{132}
\DeclareMathSymbol{\fishhookright}{\mathbin}{symbolsC}{74}

% License boxes

\newcommand*{\ollicense}{%
\begin{minipage}[b]{1.2cm}
\includegraphics[width=1.2cm]{\olpath/assets/logos/by}
\includegraphics[width=1.2cm]{\olpath/assets/logos/cc}
\end{minipage}
\hspace{.3cm}
\begin{minipage}[b]{3.5cm}
\ollicensefont
\textit{\href{https://github.com/OpenLogicProject/OpenLogic}{The Open
    Logic Text}} by the \href{http://openlogicproject.org/}{Open Logic
  Project} is licensed under a
\href{http://creativecommons.org/licenses/by/4.0/}{Creative Commons
  Attribution 4.0 International License}.
\end{minipage}
\hfill
\href{http://openlogicproject.org/}{\includegraphics[width=2.4cm]{\olpath/assets/logos/openlogic-logo-bw}}}

% - `\ifgitinfo`: produced Git revision information if `gitinfo` is used

\newcommand*{\ifgitinfo}{\ifdefined\gitVtags\sffamily\footnotesize
  rev\gitVtags:
  \href{https://github.com/OpenLogicProject/OpenLogic/commits/master}
       {\gitAbbrevHash}
  (\gitAuthorDate)\else\relax\fi}

\newcommand*{\ollicensefont}{\sffamily\fontsize{9pt}{12pt}\selectfont}

\newcommand*{\shortollicense}{\ollicensefont
  \textit{\href{https://github.com/OpenLogicProject/OpenLogic}{\jobname}}
  \ifgitinfo\ by \href{http://openlogicproject.org/}{OLP} /
  \href{http://creativecommons.org/licenses/by/4.0/}{CC--BY}}

\newcommand*{\oluselicense}[2]{%
\begin{minipage}[b]{.9cm}
\includegraphics[width=.9cm]{\olpath/assets/logos/by}
\includegraphics[width=.9cm]{\olpath/assets/logos/cc}
\includegraphics[width=.9cm]{\olpath/assets/logos/remix}
\end{minipage}
\hspace{.3cm}
\begin{minipage}[b]{5cm}
\ollicensefont
\textit{#1} by #2 is licensed under a
\href{http://creativecommons.org/licenses/by/4.0/}{Creative Commons
  Attribution 4.0 International License}. It is based on
\textit{\href{https://github.com/OpenLogicProject/OpenLogic}{The Open
    Logic Text}} by the \href{http://openlogicproject.org/}{Open Logic
  Project}, used under a
\href{http://creativecommons.org/licenses/by/4.0/}{Creative Commons
  Attribution 4.0 International License}.
\end{minipage}
\hfill
\href{http://openlogicproject.org/}{\includegraphics[width=2.7cm]{\olpath/assets/logos/openlogic-logo-bw}}}

\newcommand*{\shortoluselicense}[2]{\ollicensefont
  \textit{#1}
  by #2 based on \href{http://openlogicproject.org/}{OLP} /
  \href{http://creativecommons.org/licenses/by/4.0/}{CC--BY}}

\newcommand\gitissue[1]
  {\href{https://github.com/OpenLogicProject/OpenLogic/issues/#1}{issue
    \##1}}

\makeatother

% Photos

% Photo credits
\newlistof{photocredits}{pcr}{Photo Credits}

\newlength{\olphotowidth}
\setlength{\olphotowidth}{.5\textwidth}
\setlength{\intextsep}{0pt}
\setlength{\columnsep}{0pt}

\tagtrue{olphotos} % include photos
\tagtrue{olphotos-lowres} % pick low-res images for smaller file size

\newcommand*{\olphoto}[2]{%
  \iftag{olphotos}{%
    \IfFileExists{\olpath/assets/photos/#1/#1-large.png}{%
      \begin{wrapfigure}{R}{\olphotowidth}
        \hfil
        \iftag{olphotos-lowres}{%
          \includegraphics[width=0.8\olphotowidth]{\olpath/assets/photos/#1/#1-small}}{
          \includegraphics[width=0.8\olphotowidth]{\olpath/assets/photos/#1/#1-large}}
        \caption{#2}
      \IfFileExists{\olpath/assets/photos/#1/#1-credit.tex}{%
        \CatchFileDef{\photocredit}{\olpath/assets/photos/#1/#1-credit.tex}{}%
        \addtocontents{pcr}{#2, p.~\thepage: \photocredit\par}}{}
\end{wrapfigure}}}{}}

\newcommand*{\olasset}[1]{%
  \centerline{%
    \IfFileExists{#1}{\includegraphics[width=\olphotowidth]{#1}}
    {\includegraphics[width=\olphotowidth]{\olpath/#1}}}}


% we want all the problems deferred to the end
% OpenLogic Project
% Package for deferring processing of environments

% This is a hack using the answers package.
% A clean solution is needed

\NeedsTeXFormat{LaTeX2e}

\ProvidesFile{open-logic-defer.sty}

\RequirePackage{etoolbox}
\RequirePackage{answers}

% write problems to \oljobname.prb

\def\solutionextension{prb}

\makeatletter

% Undefine the prob and probtag environments

\csundef{prob}
\csundef{endprob}
\csundef{probtag}
\csundef{endprobtag}

% Instead, make them into environments which write probdeferred and
% probdeferredtag environments to the problem file using answers.sty

\Newassociation{prob}{probdeferred}{\jobname}
\Newassociation{probtag}{probdeferredtag}{\jobname}

% probdeferred and probdeferredtag are just wrappers for theorem
% environments which do the actual typesetting.  answers.sty writes
% the number of the problem into the .prb file, but we throw them away
% and just count them separately. It also throws away the tags. This
% is a bug; actually probtag should deal with the tags but this
% requires hacking answers.sty code.

\renewenvironment{probdeferred}[1]{\begin{probd}}{\end{probd}}
\renewenvironment{probdeferredtag}[2]{\begin{probd}}{\end{probd}}

% the probd environment does the actual typesetting.

\declaretheorem[
  style=definition, numbered=yes,
  name={Problem}]{probd}

% \probch holds the chapter number

\newcommand*{\probch}{0}

% Print the chapter number with the problem number

\renewcommand*{\theprobd}{\probch.\arabic{probd}}


% make sure the .prb file is open
\Opensolutionfile{\jobname}

% - `\problemsperchapter`: execute this in the preamble to print the
% - problem file at the end of every chapter

\newcommand{\problemsperchapter}{% 
% after each chapter command, open the .prb file and write a section
% header into it
\renewcommand\memendofchapterhook{%
  \Closesolutionfile{\jobname}
  \Opensolutionfile{\jobname}
  \Writetofile{\jobname}{
    \string\section*{Problems}
    \string\renewcommand*{\string\probch}{\thechapter}
    \string\setcounter{probd}{0}}}
  % at the end of every chapter, close the .prb file and print it
  \appto\OLEndChapterHook{%
    \Closesolutionfile{\jobname}
    \renewcommand*{\probch}{\thechapter}
    \Readsolutionfile{\jobname}}}

% - `\collectproblems`: execute this in the preamble to arrange for
% the printing of problems explicitly, e.g., at the end of the book

\newcommand{\collectproblems}{%
% at the end of every \chapter command, write a section header into
% the .prb file and reset counters
\renewcommand\memendofchapterhook{%
  \Writetofile{\jobname}{
    \string\section*{Problems for Chapter~\thechapter}
    \string\renewcommand*{\string\probch}{\thechapter}
    \string\setcounter{probd}{0}}}}

% \collectproblems is the default, so we execute it right away

\collectproblems

% - `\stopproblems`: execute at the end of the mainmatter, to stop
% chapter commands from writing to the .prb file

\newcommand{\stopproblems}{%
  \renewcommand\memendofchapterhook{}}

% - `\printproblems`: closes problem file and prints problems.

\newcommand{\printproblems}{%
  \Closesolutionfile{\jobname}
  \Readsolutionfile{\jobname}}

\makeatother



\problemsperchapter


% let's set the whole thing in Palatino, with Helvetica for
% sans-serif, and spread the lines a bit to make the text more
% readable

\usepackage{mathpazo}
\usepackage[scaled=0.95]{helvet}
\linespread{1.05}

% e-book-friendly version:
%\usepackage[papersize={3.6in,4.8in},hmargin=0.1in,vmargin={0.1in,0.1in}]{geometry}  % page geometry


\begin{document}

% First we make a titlepage

\begin{titlingpage}
\begin{raggedleft}
\fontsize{52pt}{2em}\selectfont\bfseries\sffamily
Logical\\[.5ex] 
Theory\\[.5ex] 
Text
\vskip 4ex
%\normalfont\Huge\textbf{\href{http://openlogicproject.org/}{Open Logic Project}}
%\vskip 1ex
\Large\today

\end{raggedleft}

\vfill

% oluselicense generates a license mark that a) licenses the result
% under a CC-BY licence and b) acknowledges the original source (the
% OLP).  Acknowledgment of the source is a requirement under the
% conditions of the CC-BY license used by the OLP, but you are not
% required to license the product itself under CC-BY.

\oluselicense
% Title of this version of the OLT with link to source
{\textit{Logical Theory Text}}
% Author of this version
{Fredrik Engstr\"om}
%{\href{http://openlogicproject.org/}{OLP}}
\end{titlingpage}

\frontmatter
\pagestyle{ruled}

%\newpage
\thispagestyle{empty}%
\noindent This text is a remix of the \textit{Open Logic Text} that is made for the course Logical theory, LOG110 at the University of Gothenburg by Fredrik Engstr\"om. Please see \href{http://openlogicproject.org/}{openlogicproject.org} for
more information. 

Included in this text are two sections from Tim Button's book \textit{Metatheory} which is generously realeased under a Creative Commons license making it possible to include parts of it here. Thank you!


This version of the text was compiled on \today.  
	

\tableofcontents*

\mainmatter

% olimport includes an entire part

%\olimport*[sets-functions-relations]{sets-functions-relations}

% you can also import individual chapters, but then don't forget to
% include part headings

\part{Propositional Logic}

%\olimport*[propositional-logic/{propositional-logic}

%\chapter{Syntax and Semantics}

%\olimport*[propositional-logic/syntax-and-semantics]{introduction}
%\olimport*[propositional-logic/syntax-and-semantics]{formulas}
%\olimport*[propositional-logic/syntax-and-semantics]{preliminaries}
%\olimport*[propositional-logic/syntax-and-semantics]{valuations-sat}
%\olimport*[propositional-logic/syntax-and-semantics]{semantic-notions}

%\olimport*[propositional-logic/syntax-and-semantics]{syntax-and-semantics}

\chapter{Syntax and Semantics}
\olimport*[propositional-logic/syntax-and-semantics]{introduction}
\olimport*[propositional-logic/syntax-and-semantics]{formulas}
\olimport*[propositional-logic/syntax-and-semantics]{preliminaries}
\olimport*[propositional-logic/syntax-and-semantics]{valuations-sat}
\olimport*[propositional-logic/syntax-and-semantics]{semantic-notions}
%%\documentclass[../../include/open-logic-section]{subfiles}

%\newcommand*{\define}[1]{\textsc{\lowercase{#1}}}
\newcommand*{\define}[1]{\emph{#1}}


%\begin{document}

%\olfileid{pl}{syn}{int}

\olsection{Normal forms}

%\chapter{Normal forms}%\label{ch:normalform}

In this chapter, we prove two \emph{normal form} theorems for TFL. These guarantee that, for any sentence, there is a tautologically equivalent sentence in some canonical normal form. Moreover, we shall give methods for finding these normal-form equivalents.

\subsection{Disjunctive Normal Form}%\label{s:DNFDefined}

Say that a sentence is in \define{disjunctive normal form} \emph{iff} it meets all of the following conditions:
	\begin{itemize}
		\item[(\textsc{dnf1})] No connectives occur in the sentence other than negations, conjunctions and disjunctions;
		\item[(\textsc{dnf2})] Every occurrence of negation has minimal scope (i.e.\ any `$\lnot$' is immediately followed by an atomic sentence);
		\item[(\textsc{dnf3})] No disjunction occurs within the scope of any conjunction.
	\end{itemize}
(For a reminder of the definition of the \define{scope} of a connective, see XX) So, here are are some sentences in disjunctive normal form:
	\begin{center}
		$A$\\
		$(A \land B) \lor (A \land \lnot B)$\\
		$(A \land B) \lor (A \land  B \land C \land \lnot D \land \lnot E)$\\
		$A \lor (C \land \lnot P_{234} \land P_{233} \land Q) \lor \lnot B$
	\end{center}
Note that I have here broken one of the maxims of this book (see \S\ref{s:InductionOnLength}) and \emph{temporarily} allowed myself to employ the relaxed bracketing-conventions that allow conjunctions and disjunctions to be of arbitrary length (see XXX). These conventions make it easier to see when a sentence is in disjunctive normal form. I shall continue to help myself to these relaxed conventions, without further comment, until \S\ref{s:DNFsubstitution}. 

To further illustrate the idea of disjunctive normal form, I shall introduce some more notation. I write `$\pm\meta{A}$' to indicate that $\meta{A}$ is an atomic sentence which may or may not be prefaced with an occurrence of negation. Then a sentence in disjunctive normal form has the following shape:
	$$(\pm \meta{A}_1 \land \ldots \land \pm \meta{A}_i) \lor (\pm \meta{A}_{i+1} \land \ldots \land \pm\meta{A}_j) \lor \ldots \lor (\pm\meta{A}_{m+1} \land \ldots \land \pm \meta{A}_n)$$
We now know what it is for a sentence to be in disjunctive normal form. The result that we are aiming at is:
	\begin{thm}\label{thm:dnf}\textbf{Disjunctive Normal Form Theorem.} For any sentence, there is a tautologically equivalent sentence in disjunctive normal form.
	\end{thm}\noindent 
Henceforth, I shall abbreviate `Disjunctive Normal Form' by `DNF'. 


\subsection{Proof of DNF Theorem via truth tables}\label{s:DNFTruthTable}
My first proof of the DNF Theorem employs truth tables. I shall first illustrate the technique for finding an equivalent sentence in DNF, and then turn this illustration into a rigorous proof. 

Let's suppose we have some sentence, $\meta{S}$, which contains three atomic sentences, `$A$', `$B$' and `$C$'. The very first thing to do is fill out a complete truth table for $\meta{S}$. Maybe we end up with this:
\begin{center}
\begin{tabular}{c c c | c}
$A$ & $B$ & $C$ & $\meta{S}$\\
\hline
 T & T & T & T \\
 T & T & F & F \\
 T & F & T & T \\
 T & F & F & F \\
 F & T & T & F \\
 F & T & F & F \\
 F & F & T & T \\
 F & F & F & T
\end{tabular}
\end{center}
%Now, consider a sentence, whose only connectives are negations and conjunctions, where no connective occurs within the scope of any negation, e.g.:
%	$$A \land \lnot B \land C$$
%This sentence is true when, and only when, `$A$' is true, `$B$' is false and `$C$' is true. Similarly, the sentence:
%	$$\lnot A \land \lnot B \land C$$
%this is true when, and only when, `$A$' is false, `$B$' is false and `$C$' is true. 
%
%A disjunction is true when, and only when, at least one of the disjuncts is true. So if I write down a disjunction of sentences of the above form, perhaps
%	$$(A \land \lnot B \land C) \lor (\lnot A \land \lnot B \land C)$$
%then it will be true on exactly \emph{two} lines of the truth table which describes all possible valuations of `$A$', `$B$' and `$C$'. 
%
As it happens, $\meta{S}$ is true on four lines of its truth table, namely lines 1, 3, 7 and 8. Corresponding to each of those lines, I shall write down four sentences, whose only connectives are negations and conjunctions, where every negation has minimal scope:
	\begin{itemize}
		\item[\textbullet]  `$A \land B \land C$'\hfill which is true on line 1 (and only then)
		\item[\textbullet] `$A \land \lnot B \land C$' \hfill which is true on line 3 (and only then)
		\item[\textbullet] `$\lnot A \land \lnot B \land C$' \hfill which is true on line 7 (and only then)
		\item[\textbullet] `$\lnot A \land \lnot B \land \lnot C$' \hfill which is true on line 8 (and only then)
	\end{itemize}
But if I now disjoin all of these conjunctions, like so:
$$(A \land B \land C) \lor (A \land \lnot B \land C) \lor (\lnot A \land \lnot B \land C) \lor (\lnot A \land \lnot B \land \lnot C)$$
I have a sentence in DNF which is true on exactly those lines where one of the disjuncts is true, i.e.\ it is true on (and only on) lines 1, 3, 7, and 8. So this sentence has exactly the same truth table as $\meta{S}$. So I have a sentence in DNF that is tautologically equivalent to $\meta{S}$. Which is exactly what I required!

Now, the strategy that I just adopted did not depend on the specifics of $\meta{S}$; it is perfectly general. Consequently, we can use it to obtain a simple proof of the DNF Theorem.

\begin{proof}[Proof of DNF Theorem via truth tables]
Pick any arbitrary sentence, $\meta{S}$, and let $\meta{A}_1, \ldots, \meta{A}_n$ be the atomic sentences that occur in $\meta{S}$. To obtain a sentence in DNF that is tautologically equivalent $\meta{S}$, we consider $\meta{S}$'s truth table. There are two cases to consider:
	\begin{ebullet}
		\item[\emph{Case 1: $\meta{S}$ is false on every line of its truth table.}] Then, $\meta{S}$ is a contradiction. In that case, the contradiction $(\meta{A}_1 \land \lnot \meta{A}_1) \tautequivalent \meta{S}$, and $(\meta{A}_1 \land \lnot \meta{A}_1)$ is in DNF. 
	
		\item[\emph{Case 2:  $\meta{S}$ is true on at least one line of its truth table.}]
		For each line $i$ of the truth table, let $\meta{B}_i$ be a conjunction of the form 
		$$(\pm\meta{A}_1 \land \ldots \land \pm\meta{A}_n)$$
		where the following rules determine whether or not to include a negation in front of each atomic sentence:
			\begin{align*}
				\meta{A}_m\text{ is a conjunct of }\meta{B}_i&\emph{ iff }\meta{A}_m\text{ is true on line }i\\
				\lnot\meta{A}_m\text{ is a conjunct of }\meta{B}_i&\emph{ iff }\meta{A}_m\text{ is false on line }i
			\end{align*}
		Given these rules, a trivial proof by induction shows that $\meta{B}_i$ is true on (and only on) line $i$ of the truth table which considers all possible valuations of $\meta{A}_1, \ldots, \meta{A}_n$ (i.e.\ $\meta{S}$'s truth table). 
		
		Next, let $i_1, i_2, \ldots, i_m$ be the numbers of the lines of the truth table where $\meta{S}$ is \emph{true}. Now let $\meta{D}$ be the sentence:
		$$\meta{B}_{i_1} \lor \meta{B}_{i_2} \lor \ldots \lor \meta{B}_{i_m}$$
		Since $\meta{S}$ is true on at least one line of its truth table, $\meta{D}$ is indeed well-defined; and in the limiting case where $\meta{S}$ is true on exactly one line of its truth table, $\meta{D}$ is just $\meta{B}_{i_1}$, for some $i_1$.
		
		By construction, $\meta{D}$ is in DNF. Moreover, by construction, for each line $i$ of the truth table: $\meta{S}$ is true on line $i$ of the truth table \emph{iff} one of $\meta{D}$'s disjuncts (namely, $\meta{B}_i$) is true on, and only on, line $i$. (Again, this is shown by a trivial proof by induction.) Hence $\meta{S}$ and $\meta{D}$ have the same truth table, and so are tautologically equivalent.
	\end{ebullet}
	These two cases are exhaustive and, either way, we have a sentence in DNF that is tautologically equivalent to $\meta{S}$.
\end{proof}\noindent
So we have proved the DNF Theorem. Before I say any more, though, I should immediately flag that I am hereby returning to the austere definition of a (TFL) sentence, according to which we can assume that any conjunction has exactly two conjuncts, and any disjunction has exactly two disjuncts.\footnote{NB: if you did the practice exercises for chapter \ref{ch:Induction}, you will know how to justify these conventions rigorously; and you will also, essentially, have performed the two `trivial proofs by induction' mentioned in the preceding proof.}  


\subsection{Proof of DNF Theorem via substitution}\label{s:DNFsubstitution}
I now want to offer a second proof of the DNF Theorem. 

At this point, you might reasonably ask: Why bother with \emph{two} proofs? Certainly one proof is sufficient to establish a result. But there are several possible reasons for offering new proofs of known results. For example:
	\begin{ebullet}
		\item a new proof might offer more insight (or different insights) than the original proof;
		\item a new proof-technique might be useful in other ways: it might help you to prove other results (that you have not yet proved);
		\item sheer curiosity;
	\end{ebullet}
and more besides. All these reasons apply here. 

Our first proof made use of truth tables. Our second proof will employ Lemma \ref{lem:SubstitutionTautEquivalent}. The aim is to define an \emph{algorithm} which will slowly replace subsentences with tautologically equivalent subsentences, pushing us in the direction of a DNF sentence, whilst preserving tautological equivalence, until eventually we end up with a tautologically equivalent sentence in DNF. (For a reminder about what an algorithm is, see \S\ref{s:Interpolation}.) 

That is the big idea. But the details are conceptually trickier than those in the first proof. So we shall proceed slowly. I shall start with a relatively informal overview of the algorithm, which we shall call \intoDNF{}, and shall show how it applies to a particular input sentence. With the ideas fixed by the example, I shall offer a rigorous proof that the algorithm does what we require.

Our algorithm, \intoDNF{}, takes an input sentence, $\meta{S}$, and then brings us closer and closer to DNF. The initial step of the algorithm, then, is obvious:
	\begin{ebullet}
		\item[\emph{Step 0.}] Write down your initial sentence, $\meta{S}$.
	\end{ebullet}
To make matters tractable, we shall work a particular application of the algorithm as we describe it. So, let us suppose that this is our input sentence:
	\begin{align*}
	 ((\lnot A \lif \lnot (B \lor \lnot C)) &\lif \lnot (D \liff \lnot E))
	\end{align*}
The next step of the algorithm is designed to yield a tautologically equivalent sentence which satisfies (\textsc{dnf1}). Otherwise put, our aim is to eliminate all occurrences of conditionals and biconditionals from a sentence, whilst preserving tautological equivalence. So, here is what we should do:
	\begin{ebullet}
		\item[\emph{Step 1a.}] Replace all subsentences of the form $(\meta{A} \lif \meta{B})$ with $(\lnot\meta{A} \lor \meta{B})$. 
		\item[\emph{Step 1b.}] Replace all subsentences of the form  $(\meta{A} \liff \meta{B})$ with $((\meta{A} \land \meta{B}) \lor (\lnot \meta{A} \land \lnot\meta{B}))$. 
	\end{ebullet}
Here is the result of running through Step 1 on our example. (The markers on the left-side indicate the Step we have just finished.)
{\small	\begin{align*}
	 0: & & ((\lnot A \lif \lnot (B \lor \lnot  C)) &\lif \lnot (D \liff \lnot E))\\
	 \text{1a:} & & (\lnot (\lnot \lnot A \lor \lnot (B \lor \lnot  C)) &\lor \lnot (D \liff \lnot E)) \\
	 \text{1b:} & & (\lnot (\lnot \lnot A \lor \lnot (B \lor \lnot  C)) &\lor \lnot ((D \land \lnot E) \lor (\lnot D \land \lnot \lnot E)))
	\end{align*}}
We are a little closer to DNF, since this clearly satisfies (\textsc{dnf1}). 

The next step of \intoDNF{} is designed to bring us to a tautologically equivalent sentence which satisfies (\textsc{dnf2}). Roughly, we need to `push' the negations as deeply into the sentence as we possibly can. This is what the second step of \intoDNF{} does:
	\begin{ebullet}
		\item[\emph{Step 2a:}] Replace all subsentences of the form $\lnot\lnot\meta{A}$ with $\meta{A}$.
		\item[\emph{Step 2b:}] Replace the first subsentence of the form  $\lnot (\meta{A} \land \meta{B})$ with $(\lnot \meta{A} \lor \lnot \meta{B})$.\footnote{Here `first' means `first, working from left to right'. In fact, inspecting the proof of Lemma \ref{lem:IntoDNFHalts}, it is clear that the \emph{order} in which we perform these manipulations does not matter at all; we just need to do them one at a time.}
		\item[\emph{Step 2c:}] Replace the first subsentence of the form  $\lnot (\meta{A} \lor \meta{B})$ with $(\lnot \meta{A} \land \lnot \meta{B})$. 
		\item[\emph{Step 2d:}] Repeat Steps 2a--2c, until no further change can occur.\footnote{And we can check \emph{that} no further change can occur, just by repeatedly running through Steps 2a--2c, until changes stop occurring.}
	\end{ebullet}
Applying Step 2 to to our example, we get the following sequence:
{\small	\begin{align*}
		\text{1:} & & (\lnot (\lnot \lnot A \lor \lnot (B \lor \lnot C)) &\lor \lnot ((D \land \lnot E) \lor (\lnot D \land \lnot \lnot E)))\\
		\text{2a:} & & (\lnot (A \lor \lnot (B \lor \lnot  C)) &\lor \lnot ((D \land \lnot E) \lor (\lnot D \land  E)))\\
		\text{2b:} & & (\lnot (A \lor \lnot (B \lor \lnot  C)) &\lor \lnot ((D \land \lnot E) \lor (\lnot D \land  E)))\\
		\text{2c:} & & ((\lnot A \land \lnot \lnot (B \lor \lnot  C)) &\lor \lnot ((D \land \lnot E) \lor (\lnot D \land  E)))\\
		\text{2a:} & & ((\lnot A \land (B \lor \lnot  C)) &\lor \lnot ((D \land \lnot E) \lor (\lnot D \land  E)))\\
		\text{2b:} & & ((\lnot A \land (B \lor \lnot  C)) &\lor \lnot ((D \land \lnot E) \lor (\lnot D \land  E)))\\
		\text{2c:} & & ((\lnot A \land (B \lor \lnot  C)) &\lor  (\lnot(D \land \lnot E) \land \lnot (\lnot D \land  E)))\\
		\text{2a:} & & ((\lnot A \land (B \lor \lnot  C)) &\lor  (\lnot(D \land \lnot E) \land \lnot (\lnot D \land  E)))\\
		\text{2b:} & & ((\lnot A \land (B \lor \lnot  C)) &\lor  ((\lnot D \land \lnot \lnot E) \land \lnot (\lnot D \land  E)))\\
		\text{2c:} & & ((\lnot A \land (B \lor \lnot  C)) &\lor  ((\lnot D \lor \lnot \lnot E) \land \lnot (\lnot D \land  E)))\\
		\text{2a:} & & ((\lnot A \land (B \lor \lnot  C)) &\lor  ((\lnot D \lor E) \land \lnot (\lnot D \land  E)))\\
		\text{2b:} & & ((\lnot A \land (B \lor \lnot  C)) &\lor  ((\lnot D \lor E) \land (\lnot \lnot D \lor  \lnot E)))\\
		\text{2c:} & & ((\lnot A \land (B \lor \lnot  C)) &\lor  ((\lnot D \lor E) \land (\lnot \lnot D \lor  \lnot E)))\\
		\text{2a:} & & ((\lnot A \land (B \lor \lnot  C)) &\lor  ((\lnot D \lor E) \land (D \lor  \lnot E))) 
	\end{align*}}\noindent 
Note that the sentence does not necessarily change at every sub-step. But after the last written step, no further changes \emph{can} occur; so we are done with Step 2. And we are even closer to DNF, since this satisfies (\textsc{dnf2}). 

All that remains is to meet condition (\textsc{dnf3}), i.e.\ to ensure that no disjunction occurs within the scope of any conjunction. Roughly, we need to `push' the conjunctions into the scope of disjunctions. And that is the idea behind the final step of \intoDNF:
	\begin{ebullet}
		\item[\emph{Step 3a:}]  Replace the first subsentence of the form $(\meta{A} \land (\meta{B} \lor \meta{C}))$ with $((\meta{A} \land \meta{B}) \lor (\meta{A} \land \meta{C}))$. 		
		\item[\emph{Step 3b:}]  Replace the first subsentence of the form $((\meta{A} \lor \meta{B}) \land \meta{C})$ with $((\meta{A} \land \meta{C}) \lor (\meta{B} \land \meta{C}))$. 
	\item[\emph{Step 3c:}] Repeat Steps 3a--3b, until no further changes can occur.
	\end{ebullet}
Applying Step 3 to our example, we get:
{\small	\begin{align*}
		\text{2}: & & ((\lnot A \land (B \lor \lnot  C)) &\lor  ((\lnot D \lor E) \land (D \lor  \lnot E)))\\
		\text{3a}: & & (((\lnot A \land B) \lor (\lnot A \land \lnot  C)) &\lor  ((\lnot D \lor E) \land (D \lor  \lnot E)))\\	
		\text{3b}: & & (((\lnot A \land B) \lor (\lnot A \land \lnot  C)) &\lor  ((\lnot D \land (D \lor \lnot E)) \lor (E \land (D \lor  \lnot E))))\\
		\text{3a}: & & (((\lnot A \land B) \lor (\lnot A \land \lnot  C)) &\lor  (((\lnot D \land D) \lor (\lnot D \land \lnot E)) \lor (E \land (D \lor  \lnot E))))\\
		\text{3b}: & & (((\lnot A \land B) \lor (\lnot A \land \lnot  C)) &\lor  (((\lnot D \land D) \lor (\lnot D \land \lnot E)) \lor (E \land (D \lor  \lnot E))))\\
		\text{3a}: & & (((\lnot A \land B) \lor (\lnot A \land \lnot  C)) &\lor  (((\lnot D \land D) \lor (\lnot D \land \lnot E)) \lor ((E \land D) \lor  (E \land \lnot E))))
	\end{align*}}\noindent 
And this is in DNF, as may become clearer if we \emph{temporarily} allow ourselves to remove brackets, following the notational conventions of \forallxref{} \S10.3:
		\begin{align*}
		(\lnot A \land B) \lor (\lnot A \land \lnot  C) \lor  (\lnot D \land D) \lor (\lnot D \land \lnot E) \lor (E \land D) \lor  (E \land \lnot E)
		\end{align*}
So we are done! Finally, it is worth noting that this algorithm -- although lengthy -- was \emph{much} more efficient at generating a DNF sentence than the truth-table method would have been. In fact, the truth table for our example would have $32$ lines, $22$ of which are true.

\
\\So far, I have defined the algorithm, \intoDNF, via Steps 0 through to 3c. I have also provided a worked-application of \intoDNF. But it remains to \emph{prove} that \intoDNF{} really is fit for purpose. 

First, we need to prove that any outputs from \intoDNF{} are tautologically equivalent to the inputs. This is actually quite simple:
\begin{lem}\label{lem:IntoDNFEquivalent}
	For any input sentence: any outputs of \intoDNF{} are tautologically equivalent to the input sentence.
	\begin{proof}
		At Step 0, we simply write our input. Each instance of Steps 1a--3b then involves replacing a subsentence with a tautologically equivalent subsentence. In particular, we make use of the following equivalences:
	\begin{align*}
		\text{in Step 1a:} & & (\meta{A} \lif \meta{B}) &\tautequivalent (\lnot \meta{A} \lor \meta{B})\\
		\text{in Step 1b:} & & (\meta{A} \liff \meta{B}) &\tautequivalent ((\meta{A} \land \meta{B}) \lor (\lnot \meta{A} \land \lnot\meta{B}))\\
		\text{in Step 2a:} & & \lnot \lnot \meta{A} &\tautequivalent \meta{A}\\
		\text{in Step 2b:} & & \lnot (\meta{A} \land \meta{B})&\tautequivalent(\lnot \meta{A} \lor \lnot \meta{B})\\
		\text{in Step 2c:} & & \lnot (\meta{A} \lor \meta{B}) &\tautequivalent (\lnot \meta{A} \land \lnot \meta{B})\\
		\text{in Step 3a:} & & 		(\meta{A} \land (\meta{B} \lor \meta{C})) &\tautequivalent ((\meta{A} \land \meta{B}) \lor (\meta{A} \land \meta{C}))\\
		\text{in Step 3b:} & & ((\meta{A} \lor \meta{B}) \land \meta{C}) &\tautequivalent ((\meta{A} \land \meta{C}) \lor (\meta{B} \land \meta{C}))
		\end{align*}
	Now, by Lemma \ref{lem:SubstitutionTautEquivalent}, if we replace some subsentence of a sentence with a tautologically equivalent subsentences, we obtain a sentence that is tautologically equivalent to what we started with. So any instance of any Step of \intoDNF{} preserves tautological equivalence. So, provided the algorithm yields an output within finitely many steps, that output is indeed (by a trivial induction) tautologically equivalent to the input.
	\end{proof}
\end{lem}\noindent
The next task is to show that \intoDNF{} \emph{always} yields an output. We need to ensure, for example, that the algorithm never gets trapped on an infinite loop. At the same time as we prove this, we shall also show that the output is in DNF, as we would hope. 

Unfortunately, the proof is surprisingly difficult. So I shall start with an intuitive overview of how the proof. We will aim, intuitively, to show each of the following:
\begin{ebullet}	
	\item \emph{The output of Step 1 is a sentence satisfying} (\textsc{dnf1}).

	This will be fairly trivial.
	
	\item \emph{The output of Step 2 is a sentence also satisfying} (\textsc{dnf2}). 
	
	Intuitively, this is because the negation-signs get pushed progressively deeper into the formula, until they precede only atomic sentences.
	
	\item \emph{The output of Step 3 is a sentence also satisfying} (\textsc{dnf3}).
	
	Intuitively, this is because the conjunction-signs get pushed progressively deeper into the formula, until they have no disjunctions in their scope.
\end{ebullet}
That is the main idea. As a student, your aim should be to grasp this main idea fully. Morally speaking, it just \emph{must} be right. And, for just this reason, many books at this level leave matters here, and claim that they have proved the DNF Theorem. But, \emph{just because we can},\footnote{Can, and \emph{therefore should}? There is a deep philosophical-cum-technical question: how much rigour does a proof require? But this is a book in formal logic, not philosophical logic; so I leave this as a cliffhanger.} here is a full proof, which brings these intuitive points into sharp relief:
\begin{lem}\label{lem:IntoDNFHalts}
	On any input, \intoDNF{} is guaranteed to terminate, yielding an output in DNF.
	\begin{proof}
		\emph{Concerning Step 0:} trivially, we will complete Step 0 (eventually), since this Step just involves writing down a single (finitely long) input sentence.
		
		\emph{Concerning Step 1:} The input sentence only contained a finite number of instances of `$\lif$' and `$\liff$'. These will all be eliminated (eventually) during Step 1. As such, the output at the end of Step 1 satisfies (\textsc{dnf1}). 
		
		\emph{Concerning Step 2:} This is a bit harder. For each sentence $\meta{A}$, we define its not-height, $nh(\meta{A})$, recursively:\footnote{Why do we define not-height this way? The glib answer is: because \emph{it works} in the proof. In fact, this glib answer really says it all. I have deliberately manufactured an arithmetical function, $nh$, with no other end in mind, than that successive applications of Steps 2a--2c of the algorithm are guaranteed to yield a decreasing sequence of $nh$ values.}
			\begin{align*}
				nh(\meta{A}) &= 1, \text{if }\meta{A}\text{ is atomic}\\
				nh(\lnot \meta{A}) &= 3\times{nh(\meta{A})} - 1\\
				nh(\meta{A} \land \meta{B}) = nh(\meta{A} \lor \meta{B}) &= nh(\meta{A}) + nh(\meta{B})
			\end{align*}
		I claim that, \emph{if an instance of Step 2a, 2b or 2c has any effect on the sentence, then it reduces the sentence's not-height.} In the case of Step 2a, observe that:
		\begin{align*}
		nh(\lnot\lnot\meta{A}) &= 3 \times nh(\lnot\meta{A}) -1 \\
		&= 3\times(3 \times nh(\meta{A})-1) -1 \\
		& = 9\times nh(\meta{A}) - 4
		\end{align*}
		and this is greater than $nh(\meta{A})$ for any $\meta{A}$, since a sentence's not-height must always be at least $1$. In the case of Step 2b, observe that:
	\begin{align*}
				nh(\lnot(\meta{A} \land \meta{B})) &= 3\times (nh(\meta{A} \land \meta{B})) -1 \\
				&= 3 \times(nh(\meta{A}) + nh(\meta{B})) - 1\\\
			\end{align*}
		and this is always greater than:
		\begin{align*}
		nh(\lnot\meta{A} \lor \lnot \meta{B}) &= nh(\lnot \meta{A}) + nh(\lnot\meta{B})\\
		&= 3\times nh(\meta{A}) -1 + 3 \times nh(\meta{B}) - 1\\
		&= 3\times(nh(\meta{A}) + nh(\meta{B})) - 2
		\end{align*}
		The case of Step 2c is exactly similar, and this establishes the claim. 
		
		Consequently, repeatedly running through Steps 2a--2c yields a sequence of sentences with decreasing not-heights. Since the input sentence to Step 2 (i.e.\ the output from Step 1) has some finite not-height, after a finite number of applications of Steps 2a--2c, it will be impossible to reduce the not-height any further. This brings Step 2 to an end. Trivially, the resulting sentence still satisfies (\textsc{dnf1}). And it also satisfies (\textsc{dnf2}), since the not-height of any sentence satisfying (\textsc{dnf1}) but not satisfying (\textsc{dnf2}) clearly \emph{can} be decreased by some application of Step 2a, 2b or 2c.
		
		\emph{Concerning Step 3:} This is like Step 2. For each sentence $\meta{A}$, we define its conjunction-height, $ch(\meta{A})$, recursively:
			\begin{align*}
				ch(\meta{A}) &= 2, \text{if }\meta{A}\text{ is atomic}\\
				ch(\lnot \meta{A}) &= ch(\meta{A})\\
				ch(\meta{A} \land \meta{B}) &= ch(\meta{A}) \times ch(\meta{B})\\
				ch(\meta{A} \lor \meta{B}) &=  ch(\meta{A}) + ch(\meta{B}) + 1
			\end{align*}
		I claim that, \emph{if an instance of Step 3a or 3b has any effect on the sentence, then it reduces the sentence's conjunction-height}. In the case of Step 3a, observe that 
			\begin{align*}
				ch(\meta{A} \land (\meta{B} \lor \meta{C})) &= ch(\meta{A}) \times  (ch(\meta{B}) + ch(\meta{C}) + 1)\\
					&= ch(\meta{A})\times ch(\meta{B}) + ch(\meta{A}) \times ch(\meta{C}) + ch(\meta{A})
			\end{align*}
		and, since $ch(\meta{A}) > 1$ for any $\meta{A}$, this must be larger than:
			\begin{align*}
				ch((\meta{A} \land \meta{B}) \lor (\meta{A} \land \meta{C})) &= ch(\meta{A}) \times ch(\meta{B}) + ch(\meta{A}) \times ch(\meta{C}) + 1 
			\end{align*}
		The case of Step 3b is exactly similar. This establishes the claim. 
		
		Now, as above: running through Steps 3a--3b repeatedly yields a sequence of sentences with decreasing conjunction-heights. So, after a finite number of applications of Steps 3a--3b, it will be impossible to reduce the conjunction-height any further. Trivially, the resulting sentence still satisfies (\textsc{dnf1}) and (\textsc{dnf2}). And it also satisfies (\textsc{dnf3}), since the conjunction-height of any sentence satisfying both (\textsc{dnf1}) and (\textsc{dnf2}) but not (\textsc{dnf3}) \emph{can} be decreased by some application of Step 3a or 3b.
	\end{proof}
\end{lem}\noindent
Now if we just combine combine Lemma \ref{lem:IntoDNFEquivalent} with Lemma \ref{lem:IntoDNFHalts}, we have a new proof of the DNF Theorem, via substitution.





\subsection{Cleaning up DNF sentences}\label{s:DNFCleanUp}
I have offered two different proofs of the DNF Theorem. They also described two different algorithms for putting sentences into DNF. But these algorithms do not always leave you with the most elegant output. Indeed, just as we can `clean up' the output of the algorithm for calculating interpolants (see \S\ref{s:Interpolation}), we can `clean up' the output of our algorithms for calculating DNF-equivalents. Here are the clean-up rules; I leave it as an exercise to \emph{justify} them. First, we rewrite our DNF sentence using the relaxed notational conventions of \forallxref{} \S10.3. Now we apply the following:

\
\\\emph{Rule 1: remove repetitious conjuncts.}  If any disjunct contains the same conjunct more than once, remove all but one instance of that conjunct (respecting bracketing conventions as appropriate).
\\\emph{Example:} if the output is:
$$(A \land B) \lor (A \land \lnot C) \lor D \lor (A \land B)$$
then applying Rule 1 we get:
$$(A \land B) \lor (A \land \lnot C) \lor D$$
\noindent \emph{Rule 2: remove overly-specific disjuncts.} Delete any disjunct which entails any other disjunct. (If two or more disjuncts entail each other, delete all but the first of those disjuncts.)
\\\emph{Example:} the output of our example in \S\ref{s:DNFsubstitution} was:
		\begin{align*}
		(A \land B) \lor (A \land \lnot C) &\lor  (\lnot D \land D) \lor (\lnot E \land D) \lor (\lnot D \land E) \lor (\lnot E \land  E)		
		\end{align*}
		Noting that a contradiction entails \emph{everything}, we should remove both contradictory disjuncts, obtaining (still with relaxed conventions):
		\begin{align*}
		(A \land B) \lor (A \land \lnot C) &\lor  (\lnot E \land D) \lor (\lnot D \land E)
		\end{align*}		
\emph{Rule 3: invoke excluded middle.} Suppose there is a disjunct whose conjuncts are just $\meta{A}_1, \ldots,  \meta{A}_m$ and $\meta{B}$ (in any order), and another disjunct whose conjuncts are just $\meta{A}_1, \ldots, \meta{A}_m$ and $\lnot \meta{B}$ (in any order), so that the two disjuncts disagree only on whether whether $\meta{B}$ should be prefixed with a negation. In that case, delete both disjuncts and replace them with the simpler disjunct $(\meta{A}_1 \land \ldots \land \meta{A}_m)$. (However, if we have \emph{just} $\meta{B}$ and $\lnot\meta{B}$ as disjuncts, then delete the entire sentence and leave the tautology $(\meta{B} \lor \lnot \meta{B})$.)
\\\emph{Example:} the output of the example in \S\ref{s:DNFTruthTable} was:
		$$(A \land B \land C) \lor (A \land \lnot B \land C) \lor (\lnot A \land \lnot B \land C) \lor (\lnot A \land \lnot B \land \lnot C)$$
			which can be simplified to:
		$$(A \land C) \lor (\lnot A \land \lnot B)$$


\subsection{Conjunctive Normal Form}\label{s:CNF}
So far in this chapter, I have discussed \emph{disjunctive} normal form. Given the duality of disjunction and conjunction (see \S\ref{s:Duality}), it may not come as a surprise to hear that there is also such a thing as \emph{conjunctive normal form} (CNF).

The definition of CNF is exactly analogous to the definition of DNF. So, a sentence is in CNF \emph{iff} it meets all of the following conditions:
	\begin{itemize}
		\item[(\textsc{cnf1})] No connectives occur in the sentence other than negations, conjunctions and disjunctions;
		\item[(\textsc{cnf2})] Every occurrence of negation has minimal scope;
		\item[(\textsc{cnf3})] No conjunction occurs within the scope of any disjunction. 
	\end{itemize}
Generally, then, a sentence in CNF looks like this
	$$(\pm \meta{A}_1 \lor \ldots \lor \pm \meta{A}_i) \land (\pm \meta{A}_{i+1} \lor \ldots \lor \pm\meta{A}_j) \land \ldots \land (\pm\meta{A}_{m+1} \lor\ldots \lor \pm \meta{A}_n)$$
where each $\meta{A}_k$ is an atomic sentence.

Since `$\lnot$' is its own dual, and `$\lor$' and `$\land$' are the duals of each other, it is immediate clear that if a sentence is in DNF, then its dual is in CNF; and \emph{vice versa}. Armed with this insight, we can immediately prove another normal form theorem:
	\begin{thm}\label{thm:cnf}\textbf{Conjunctive Normal Form Theorem.} For any sentence, there is a tautologically equivalent sentence in conjunctive normal form.
\begin{proof}
	Let $\meta{S}$ be any sentence. Applying the DNF Theorem to $\dual{\meta{S}}$, there is a DNF sentence, $\meta{D}$, such that $\dual{\meta{S}} \tautequivalent \meta{D}$. So, by Theorem \ref{thm:DualEquivalence}, $\dual{\dual{\meta{S}}} \tautequivalent \dual{\meta{D}}$. Whence, by Lemma \ref{lem:DualDual}, $\meta{S} \tautequivalent \dual{\meta{D}}$. But, since $\meta{D}$ is in DNF, $\dual{\meta{D}}$ is in CNF. 	
\end{proof}
\end{thm}\noindent
This slick proof is a further illustration of the power of duality. However, it might suggest that the DNF Theorem enjoys some kind of `precedence' over the CNF Theorem. That would be misleading.  We can easily prove the CNF Theorem directly, using either of the two proof techniques that we used to prove the DNF Theorem (whereupon the DNF Theorem could be proved as a consequence of the CNF Theorem and duality). I shall sketch the main ideas of the direct proofs of the CNF Theorem; I leave it as an exercise to make these ideas more precise.
\begin{proof}[Proof sketch of CNF Theorem via truth tables]
	Given a TFL sentence, $\meta{S}$, we begin by writing down the complete truth table for $\meta{S}$.
	
	If $\meta{S}$ is \emph{true} on every line of the truth table, then $\meta{S} \tautequivalent (\meta{A}_1 \lor \lnot \meta{A}_1)$.
	
	If $\meta{S}$ is \emph{false} on at least one line of the truth table then, for every line on the truth table where $\meta{S}$ is false, write down a disjunction $(\pm\meta{A}_1 \lor \ldots \lor \pm\meta{A}_n)$ which is \emph{false} on (and only on) that line. Let $\meta{C}$ be the conjunction of all of these disjuncts; by construction, $\meta{C}$ is in CNF and $\meta{S} \tautequivalent \meta{C}$.
\end{proof}
\begin{proof}[Proof sketch of CNF Theorem via substitution.] Using Steps 0--2 of \intoDNF, we  obtain a sentence satisfying (\textsc{cnf1}) and (\textsc{cnf2}). To turn this into a sentence in CNF, we simply modify Step 3, using substitutions based on the duals (what else!)\ of the Distribution Laws used in Step 3 of \intoDNF{}; i.e.:
	\begin{align*}
		((\meta{A} \land \meta{B}) \lor \meta{C}) &\tautequivalent ((\meta{A} \lor \meta{C}) \land (\meta{B} \lor \meta{C}))\\
		(\meta{A} \lor (\meta{B} \land \meta{C})) &\tautequivalent ((\meta{A} \lor \meta{B}) \land (\meta{A} \lor \meta{C}))	
	\end{align*}
This allows us to pull conjunctions outside the scope of any disjunction.
\end{proof}\noindent
Of course, after obtaining a sentence in CNF -- however we chose to do it -- we can clean-up the resulting sentence using rules similar to those of \S\ref{s:DNFCleanUp}. I leave it as an exercise, to determine what the appropriate rules should be.

\practiceproblems
\problempart
\label{pr.DNF}
Consider the following sentences:
	\begin{itemize}
		\item $(A \lif \lnot B)$
		\item $\lnot (A \liff B)$
		\item $(\lnot A \lor \lnot (A \land B))$
		\item $(\lnot (A \lif B ) \land (A \lif C))$
		\item $(\lnot (A \lor B) \liff ((\lnot C \land \lnot A) \lif \lnot B))$
		\item $((\lnot (A \land \lnot B) \lif C) \land \lnot (A \land D))$
	\end{itemize}
For each sentence:
	\begin{ebullet}
		\item use both algorithms (by truth table, and by substitution) to write down sentences in DNF that are tautologically equivalent to these sentences.
		\item use both algorithms to write down sentences in CNF that are tautologically equivalent to these sentences. 
	\end{ebullet}

\problempart
Offer proofs of the two `trivial inductions' mentioned in the proof of Theorem \ref{thm:dnf} via truth tables.\\

\problempart
Explain why the proof of Theorem \ref{lem:IntoDNFEquivalent} required `a trivial induction' (in the last sentence). Prove that induction.\\


\problempart
Justify each of the rules mentioned in \S\ref{s:DNFCleanUp}, i.e.\ prove that applying each of them preserves tautological equivalence.\\

\problempart
Fill out the details of the two proof sketches of the CNF Theorem. Determine the appropriate `clean up' rules for CNF sentences.

%\end{document}
\OLEndChapterHook

\tagfalse{FOL}

\chapter{Natural Deduction}
\olimport*[first-order-logic/proof-systems]{introduction}
\olimport*[first-order-logic/proof-systems]{natural-deduction}
\olimport*[first-order-logic/natural-deduction]{rules-and-proofs}
\olimport*[first-order-logic/natural-deduction]{propositional-rules}
\olimport*[first-order-logic/natural-deduction]{derivations}
\olimport*[first-order-logic/natural-deduction]{proving-things}
\olimport*[first-order-logic/natural-deduction]{proof-theoretic-notions}
\olimport*[first-order-logic/natural-deduction]{provability-consistency}
\olimport*[first-order-logic/natural-deduction]{provability-propositional}
\olimport*[first-order-logic/natural-deduction]{soundness}


%\olimport*[first-order-logic/natural-deduction]{natural-deduction}
\olimport*[first-order-logic/completeness]{completeness}
\tagtrue{FOL}

\OLEndChapterHook


\part{First-order Logic}

\olimport*[first-order-logic/syntax-and-semantics]{syntax-and-semantics}

\olimport*[first-order-logic/models-theories]{models-theories}

% For a proof system, we'll do only natural deduction. But some of the
% texts will refer to other proof systems unless you set some tags in
% the config.sty file.  So make sure those are set. We'll include
% these sections by hand so we can add a couple of sections from the
% proof-systems chapter. Of course, if new sections are added or
% sections are moved or renamed in the main repository, this may break.

\chapter{Natural Deduction}

\olimport*[first-order-logic/proof-systems]{introduction}

\olimport*[first-order-logic/proof-systems]{natural-deduction}

\olimport*[first-order-logic/natural-deduction]{rules-and-proofs}

\olimport*[first-order-logic/natural-deduction]{propositional-rules}

% We'll reorder things: let's do propositional examples first and then
% go back to quantifiers

\olimport*[first-order-logic/natural-deduction]{derivations}

\olimport*[first-order-logic/natural-deduction]{proving-things}

\olimport*[first-order-logic/natural-deduction]{quantifier-rules}

\olimport*[first-order-logic/natural-deduction]{proving-things-quant}

\olimport*[first-order-logic/natural-deduction]{proof-theoretic-notions}

\olimport*[first-order-logic/natural-deduction]{provability-consistency}

\olimport*[first-order-logic/natural-deduction]{provability-propositional}

\olimport*[first-order-logic/natural-deduction]{provability-quantifiers}

\olimport*[first-order-logic/natural-deduction]{soundness}

\olimport*[first-order-logic/natural-deduction]{identity}

\olimport*[first-order-logic/natural-deduction]{soundness-identity}

% Chapters should end with \OLEndChapterHook. They will automatically
% if you include entire parts or chapters. Here we did a chapter ``by
% hand'' so we should add \OLEndChapterHook by hand too

\OLEndChapterHook

\olimport*[first-order-logic/completeness]{completeness}

\olimport*[first-order-logic/beyond]{beyond}

% OLEndPartHook should come at the end of each
% part.

\OLEndPartHook

% Include some more chapters and parts

\olimport*[turing-machines]{turing-machines}

% Part: incompleteness

\part{Computability and Incompleteness}

% include intro to recursive function from computability part

\olimport*[computability/recursive-functions]{recursive-functions}

\olimport*[incompleteness/arithmetization-syntax]{arithmetization-syntax}

\olimport*[incompleteness/representability-in-q]{representability-in-q}

% leave out this part -- it depends on computability theory chapter
% \olimport[incompleteness/theories-computability]{theories-computability}

\olimport*[incompleteness/incompleteness-provability]{incompleteness-provability}

\OLEndPartHook

\stopproblems

% Ok, that's it. Now for the appendices

\appendix

\olimport*[methods/proofs]{proofs}
\olimport*[methods/induction]{induction}
\olimport*[history/biographies]{biographies}

%\olimport*[methods]{methods}



%\olimport*[history]{history}


\backmatter

% If you include any chapters from the history part, you have to print
% the Photo Credits. 

\photocredits

% Include the bibliography

\bibliographystyle{\olpath/bib/natbib-oup}
\bibliography{\olpath/bib/open-logic}

\end{document}

